\documentclass[openany]{book}
\usepackage{epigraph}

\usepackage[margin=1in]{geometry} 
\usepackage{amsmath,amsthm,amssymb}
\usepackage{hyperref}
\hypersetup{pdfborder=0 0 0}
\usepackage{hyperref}

\usepackage{parskip}
\usepackage{graphicx}
\usepackage{algorithm}
\usepackage[noend]{algpseudocode}
 \usepackage{hyperref}
\hypersetup{
    colorlinks=true,
    linkcolor=blue,
    filecolor=magenta,      
    urlcolor=blue,
    pdftitle={Overleaf Example},
    pdfpagemode=FullScreen,
    }
\newcommand{\N}{\mathbb{N}}
\newcommand{\Z}{\mathbb{Z}}
 
\newenvironment{theorem}[2][Theorem]{\begin{trivlist}
\item[\hskip \labelsep {\bfseries #1}\hskip \labelsep {\bfseries #2.}]}{\end{trivlist}}
\newenvironment{lemma}[2][Lemma]{\begin{trivlist}
\item[\hskip \labelsep {\bfseries #1}\hskip \labelsep {\bfseries #2.}]}{\end{trivlist}}
\newenvironment{exercise}[2][Exercise]{\begin{trivlist}
\item[\hskip \labelsep {\bfseries #1}\hskip \labelsep {\bfseries #2.}]}{\end{trivlist}}
\newenvironment{reflection}[2][Reflection]{\begin{trivlist}
\item[\hskip \labelsep {\bfseries #1}\hskip \labelsep {\bfseries #2.}]}{\end{trivlist}}
\newenvironment{proposition}[2][Proposition]{\begin{trivlist}
\item[\hskip \labelsep {\bfseries #1}\hskip \labelsep {\bfseries #2.}]}{\end{trivlist}}
\newenvironment{corollary}[2][Corollary]{\begin{trivlist}
\item[\hskip \labelsep {\bfseries #1}\hskip \labelsep {\bfseries #2.}]}{\end{trivlist}}




\begin{document}


\title{Theoratical Minimum  - An ML guide for Practitioners and ML Interviews} 

\author{Compiled from various sources across the internet} 
% --------------------------------------------------------------
%                         Start here
% --------------------------------------------------------------
 
%\renewcommand{\qedsymbol}{\filledbox}
 


\maketitle

\tableofcontents

\newpage

\chapter{Introduction}


\epigraph{“All models are wrong, but some are useful.”
— George Box}






\section{What is Machine learning and its types? }


Machine learning is a subset of artificial intelligence(AI), and is evolving to be critical for all fields of AI.  Machine Learning refers to a broad set of computer science techniques that let us give computers, as Arthur Samuel put it in 1959, “the ability to learn without being explicitly programmed.” There are many different types of machine learning algorithms, including reinforcement learning, genetic algorithms, rule-based machine learning, learning classifier systems, and decision trees. The Wikipedia article has many more examples. 

\textit{Over the past decades computers have broadly automated tasks that programmers could describe with clear rules and algorithms. Modern machine learning techniques now allow us to do the same for tasks where describing the precise rules is much harder.} \\


We can think of ML in various ways. e.g ML is about 

\begin{itemize}
    \item Models that capture an underlying process  
    \item Function Approximation or Sequence Modeling (statistical view is to  find a model that fits a data set well)
    \item Programs that learn from data  
    \item A tool for Data Compression
    \item In machine learning, we often fit models, but as a means to the end of making good predictions or decisions.\\
\end{itemize}

\noindent
\textbf{Formal Definition:} 

we can formalize the notion using the definition proposed by Murphy. 

\emph{"A computer program is said to learn from experience $E$ with respect to some class of tasks $T$ and performance measure $P$, if its performance at tasks in $T$, as measured by $P$, improves with experience $E . "$} \\

A lot of value that is being created these days by AI Startups is by this technology we call supervised machine learning. Its involves showing a bunch of examples to the computer so that it can learn to generalise a certain concept from those examples. Machine Learning a subfield of AI and has been around for decades. How can pick real world problems and try to solve them using this technology? The most generic approach is to think about the ML definition. which is” A computer program is said to learn from experience E with respect to some class of tasks T and performance measure P, if its performance at tasks in T, as measured by P, improves with experience E.”

Start with the informal problem description which describes the general goals. E.g  I need a program that will tell me which tweets will get retweets.Convert it into a formal description For example:

\begin{itemize}
    \item \textbf{Task (T):} Classify a tweet that has not been published as going to get retweets or not.

    \item \textbf{Experience (E):} A corpus of tweets for an account where some have retweets and some do not.

    \item \textbf{Performance (P):} Classification accuracy, the number of tweets predicted correctly out of all tweets considered as a percentage.

\end{itemize}

Machine learning is about learning from observed data and making predictions and/or decisions. We use the observations to learn a function that can then be used for predictions and inference. Machine Learning uses a set of observations to uncover an underlying process. In Statistics, underlying process is a probability distribution. In tackling machine learning (and computer science in general) we face some deep philosophical questions. Questions like, “What does it mean to learn?” and, “Can a computer learn?” and, “How do you define simplicity?” and, “Why does Occam’s Razor work? (Why do simple hypotheses do well at modelling reality?)” In a very deep sense, learning theorists take these philosophical questions — or at least aspects of them — give them fleshy mathematical bodies, and then answer them with theorems and proofs. These fleshy bodies might have imperfections or they might only address one small part of a big question, but the more we think about them the closer we get to robust answers. 

%For a lot of companies, it can be a hypothesis. e.g Algorithm/feature/design X will increase member engagement with our service and ultimately member retention.




%Machine learning is based on \textit{probability theory and statistics and optimization} , is the basis for big data, data science, \textit{predictive modeling}, \textit{data mining, information retrieval}, etc, and becomes a critical ingredient for computer vision, natural language processing, robotics, etc. Reinforcement learning is akin to \textit{optimal control}, and operations research and management, and is also related to psychology and neuroscience.   




\section{Types of ML Systems}

Usually, we categorize machine learning as \textbf{supervised, unsupervised, and reinforcement learning.} In supervised learning, there are \textbf{labelled data}; in unsupervised learning, there are no labeled data; and in reinforcement learning, there are \textbf{evaluative feedbacks}, but no supervised signals.  Classification and regression are two types of supervised learning problems, with categorical and numerical outputs respectively.

A supervised machine learning algorithm is composed of \textit{a dataset, a cost/loss function, an optimization procedure, and a model}.  A dataset is divided into non-overlapping training, validation,  and testing subsets.  A cost/loss function measures the model performance,  e.g.,  with respect to the accuracy, like mean square error in regression and classification error rate. Training error measures the error on the training data, minimizing which is an optimization problem.  Generalization error, or test error, measures the error on new input data, which differentiates machine learning from optimization.  A machine learning algorithm tries to make the training error, and the gap between training error and testing error small.  A model is \textit{underfitting} if it can not achieve a low training error; a model \textit{is over-fitting} if the gap between training error and test error is large. \\

Example: Let's we have a dataset consisting of two columns, GDP of a country and its life satisfaction index. We can using linear regression model life satisfaction as a linear function of GDP per capital. 

%Life_satisfaction = Theta_0 + Theta_1*GDP_per_capita

%Training result;
%theta_1 = 4.85 
%theta_2 = 4.91 * 10^-5

Unsupervised learning attempts to extract information from data without labels, e.g., clustering and density estimation.  Representation learning is a classical type of unsupervised learning.  Representation learning finds a representation to preserve as much information about the original data as possible, at the same time, to keep the representation simpler  or more accessible than the original data, with low-dimensional, sparse, and independent representations.

 Reinforcement learning is akin to \textit{optimal control}, and operations research and management, and is also related to psychology and neuroscience.   

The current darling of these machine learning algorithms is deep learning algorithms, which we'll discuss in detail in a separate guide.  \textit{Deep learning}, or deep neural networks, is a particular machine learning scheme, usually for supervised or unsupervised learning, and can be integrated with reinforcement learning, usually as a \textit{function approximator}. Supervised and unsupervised learning are usually one-shot, myopic, considering instant reward; while reinforcement learning is sequential, far-sighted, considering long-term accumulative reward. 

%However, training feedforward networks or convolutional neural networks with supervised learning is a kind of representation learning.




%ML is helpful when In unsupervised learning we use a high-level representation of the input.


\textbf{Generalization}

The ultimate goal of Machine learning is to be able to generalize to new unseen data. i.e predicting behavior under new conditions. we study Problems like: Regression, Clustering, Classification, Recommendation.

\textit{``The ability to perform well on previously unobserved inputs is called generalization".
}
How hard is this objective? well generally speaking the more data we have, the more better chance we have of figuring out the ground truth function 'f'. Our success also depends on no. of functions that could be ground truth(fewer the better). But again, Machine learning is not just about optimization where we reduce the training error. what we really care about is the generalization error(or test error), defined as "the expected value of the error on a new input. i.e probability that g disagrees with the ground truth f on the label of a random sample" 

Our Major goal of machine learning is to generalize beyond the examples seen during training. Memorisation is not a good strategy for learners because we often see only a fraction of examples we can will encounter in real world settings. Luckily induction works pretty well, which means by taking certain assumptions ML algorithms are able to figure out knowledge from a small number of examples.

\section{Framing ML Problems}

Machine learning enables us to tackle tasks that are too difficult to solve with fixed programs written and designed by human beings. From a scientific and philosophical point of view, machine learning is interesting because developing our understanding of it entails developing our understanding of the principles that underlie intelligence.

In this relatively formal definition of the word "task," the process of learning itself is not the task. Learning is our means of attaining the ability to perform the task. For example, if we want a robot to be able to walk, then walking is the task. We could program the robot to learn to walk, or we could attempt to directly write a program that specifies how to walk manually.

Machine learning tasks are usually described in terms of how the machine learning system should process an example. An example is a collection of features that have been quantitatively measured from some object or event that we want the machine learning system to process. We typically represent an example as a vector $\boldsymbol{x} \in \mathbb{R}^n$ where each entry $x_i$ of the vector is another feature. For example, the features of an image are usually the values of the pixels in the image.

 discriminative analysis helps us answer two important kinds of questions:
 
\begin{tabular}{|l|l|l|l|}
\hline Categorize & What is it? & aka classify & e.g. benign vs malignant? which species? \\
\hline Quantify & How much? & aka regress & e.g. price? distance? age? radioactivity? \\
\hline
\end{tabular}

Digging one level deeper, there are two types of categorization:

\begin{tabular}{|l|l|l|}
\hline Binary & Checking for the presence of a single condition & e.g. tumor \\
\hline Multi-Label & When there are many possible outcomes & e.g. species \\
\hline
\end{tabular}


To decide which ML model to use to solve your problem, you first need to frame your problem into a problem that ML can solve.

The most general types of ML problems are classification and regression. Classification models classify inputs into different categories. For example, you want to classify each email to be either spam or not spam. Regression models output a continuous value. An example is a house prediction model that outputs the price of a given house.

A regression model can easily be framed as a classification model and vice versa. For example, house prediction can become a classification task if we quantize the house prices into buckets such as under $100,000, $100,000 - 200,000, $200,000 - 500,000, … and predict the bucket the house should be in. The email classification model can become a regression model if we make it output values between 0 and 1, and use a threshold to determine which values should be SPAM (for example, if the value is above 0.5, the email is spam).

Within classification problems, the fewer classes there are to classify, the simpler the problem is. The simplest is binary classification where there are only two possible classes. Examples of binary classification include classifying whether a comment is toxic or not toxic, whether a lung scan shows signs of cancer or not, whether a transaction is fraudulent or not. It’s unclear whether this type of problem is common because they are common in nature or simply because ML practitioners are most comfortable handling them. Dealing with binary classification problems, including upsampling and downsampling as well as visualizing ROC curves and confusion matrices, is much easier than dealing with multiclass classifiers.

When there are more than two classes, the problem becomes multiclass classification. When the number of classes is high, such as disease diagnosis where the number of diseases can go up to thousands or product classifications where the number of products can go up to tens of thousands, the problem can be very challenging. The first challenge is in data collection. In my experience, ML models typically need at least 100 examples for each class to learn to classify that class. So if you have 1000 classes, you already need at least 100,000 examples. The data collection can be especially challenging when some of the classes are rare, and if you have thousands of classes, there’s a high likelihood that some of them are rare.

When the number of classes is large, hierarchical classification might be useful. In hierarchical classification, you have a classifier to first classify each example into one of the large groups. Then you have another classifier to classify this example into one of the subgroups. For example, for product classification, you can first classify each product into one of the four main categories: electronics, home & kitchen, fashion, and pet supplies. After a product has been classified into a category, say fashion, you can use another classifier to put this product into one of the subgroups: shoes, shirt, jeans, accessories.

In both binary and multiclass classification, each example belongs to exactly one class. When an example can belong to multiple classes, we have a multilabel classification problem. For example, when building a model to classify articles into three topics: [tech, entertainment, finance], an article can be in both tech and finance.

There are two major approaches to multilabel classification problems. The first is to treat it as you would a multiclass classification. In multiclass classification, if there are three possible classes [tech, entertainment, finance] and the label for an example is entertainment, you represent this label with the vector [0, 1, 0]. In multilabel classification, if an example has both labels entertainment and finance, its label will be represented as [0, 1, 1].

The second approach is to turn it into multiple binary classification problems. For the article classification problem above, you can have three models corresponding to three topics, where each model outputs whether an article is in that topic or not.

Changing the way you frame your problem might make your problem significantly harder or easier. Consider the task of predicting what app a phone user wants to use next. A naive setup would be to frame this as a multiclass classification task — use the user’s and environment’s features (user demographic information, time, location, previous apps used) as input and output a probability distribution for every single app on the user’s phone. This means that the last layer of your model will have the shape |number of hidden layer| x |number of apps|.

This is a bad approach because whenever a new app is added, you have to retrain your model, or at least the last layer of your model. A better approach is to frame this as a regression task — the user’s, the environment’s, and the app’s features as input, and output a value between 0 and 1, the higher the value, the more likely the user will open the app given the context.







\textbf{Structured output:} Structured output tasks involve any task where the output is a vector (or other data structure containing multiple values) with important relationships between the different elements. This is a broad category and subsumes the transcription and translation tasks described above, as well as many other tasks. One example is parsing-mapping a natural language sentence into a tree that describes its grammatical structure by tagging nodes of the trees as being verbs, nouns, adverbs, and so on.

\textbf{Anomaly detection:} In this type of task, the computer program sifts through a set of events or objects and flags some of them as being unusual or atypical. An example of an anomaly detection task is credit card fraud detection. By modeling your purchasing habits, a credit card company can detect misuse of your cards. If a thief steals your credit card or credit card.

\begin{tabular}{|l|l|l|}
\hline Binary & Checking for the presence of a single condition & e.g.tumor \\
\hline Multi-Label & When there are many possible outcomes & e.g. species \\
\hline
\end{tabular}


Few practical use-cases:

\begin{itemize}
    \item Object recognition
    \item Speech recognition / sound detection
    \item Natural Language Processing / Sentiment analysis
    \item Creative (e.g. Style Transfer - learning to draw an image
in the style of an artist)
    \item Prediction - given some inputs, what is the expected output for unseen examples
    \item Translation between languages
    \item Restoration / Transformation - eg taking an image and

\end{itemize}


\section{Is ML the right tool for your problem?}

Here are some more General Heuristics for picking problems that  are solvable by machine learning:

\begin{itemize}
    \item \textbf{You cannot code the rules:} Many human tasks (such as recognizing whether an email is spam or not spam) cannot be adequately solved using a simple (deterministic), rule-based solution. A large number of factors could influence the answer. When rules depend on too many factors and many of these rules overlap or need to be tuned very finely, it soon becomes difficult for a human to accurately code the rules. You can use ML to effectively solve this problem.

    \item \textbf{Complex}: the patterns are complex: Consider a website like Airbnb with a lot of house listings, each listing comes with a zip code. If you want to sort listings into the states they are located in, you wouldn’t need an ML system. Since the pattern is simple—each zip code corresponds to a known state—you can just use a lookup table. The relationship between a rental price and all its characteristics follows a much more complex pattern which would be very challenging to explicitly state by hand. ML is a good solution for this. Instead of telling your system how to calculate the price from a list of characteristics, you can provide prices and characteristics, and let your ML system figure out the pattern.


    \item \textbf{Patterns}: there are patterns to learn ML solutions are only useful when there are patterns to learn. Sane people don’t invest money into building an ML system to predict the next outcome of a fair die because there’s no pattern in how these outcomes are generated.
    However, there are patterns in how stocks are priced, and therefore companies have invested billions of dollars in building ML systems to learn those patterns.
    Whether a pattern exists might not be obvious, or if patterns exist, your dataset might not be sufficient to capture them. For example, there might be a pattern in how Elon Musk’s tweets affect Bitcoin prices. However, you wouldn’t know until you’ve rigorously trained and evaluated your ML models on his tweets. Even if all your models fail to make reasonable predictions of Bitcoin prices, it doesn’t mean there’s no pattern.

    
    \item \textbf{You cannot scale}: You might be able to manually recognized a few hundred emails and decide whether they are spam or not. However, this task becomes tedious for millions of emails. ML solutions are effective at handling large-scale problems. 

    \item \textbf{Personalization:} The application requires that the software customize to its operational environment after it is fielded. One example of this is speech recognition systems that customize to the user who purchases the software. Machine learning here provides the mechanism for adaptation. Software applications that customize to users are growing rapidly - e.g., bookstores that customize to your purchasing preferences, or email readers that customize to your particular definition of spam. This machine learning niche within the software world is growing rapidly.

    \item \textbf{Dataset availability:} You have access to a sizeable set of data from which to train your model. There are no absolutes about how much data is enough, but every feature (data attribute) that you include in your model increases the number of instances (data records) you'll need to properly train it. Also account for splitting your dataset into three subsets: one for training, one for evaluation, and one for testing. In the zero-shot learning (sometimes known as zero-data learning) context, it’s possible for an ML system to make correct predictions for a task without having been trained on data for that task. However, this ML system was previously trained on data for a related task. So even though the system doesn’t require data for the task at hand to learn from, it still requires data to learn.
    
    It’s also possible to launch an ML system without data. For example, in the context of online learning, ML models can be deployed without having been trained on any data, but they will learn from data in production. However, serving insufficiently trained models to users comes with certain risks, such as poor customer experience.
    
    Without data and without online learning, many companies follow a ‘fake-it-til-you make it’ approach: launching a product that serves predictions made by humans, instead of ML algorithms, with the hope of using the generated data to train ML algorithms.


    \item You can't use an easier and more concrete mathematical model to solve your problem.
 
\end{itemize}

%add more from here

\section{General Steps}

While we’ll take a deeper dive into what each of these steps mean in practice in later chapters, let’s take a brief look at what happens during each of the steps. 


\begin{itemize}
    \item \textbf{Step 1. Project scoping: } A project starts with scoping the project, laying out goals, objectives, and constraints. Stakeholders should be identified and involved. Resources should be estimated and allocated. We already discussed different stakeholders and some of the focuses for ML projects in production earlier in this chapter. We’ll discuss how to scope an ML project in the context of a business as well as how to organize teams to ensure the success of an ML project in Chapter 9: The Human Side of Machine Learning.

    \item \textbf{Step 2. Data engineering} A vast majority of ML models today learn from data, so developing ML models starts with engineering data. In chapter 2, we’ll discuss the importance of data in ML and the fundamentals of data engineering, which covers handling data from different sources and formats. With access to raw data, we’ll want to curate training data out of it by sampling and generating labels.

    \item \textbf{Step 3. ML model development}with the initial set of training data, we’ll need to extract features and develop initial models leveraging these features. This is the stage that requires the most ML knowledge and is most often covered in ML courses. In chapter 4, we’ll discuss feature engineering and in chapter 5, we’ll discuss model selection, training, and evaluation.

    \item \textbf{Step 4. Deployment} After a model is developed, it needs to be made accessible to users. Developing an ML system is like writing — you will never reach the point when your system is done. But you do reach the point when you have to put your system out there. We’ll discuss different ways to deploy an ML model in chapter 6.

    \item \textbf{Step 5. Monitoring and continual learning} Once in production, models need to be monitored for performance decay and maintained to be adaptive to changing environments and changing requirements. This step will be discussed in chapter 7.

    \item \textbf{Step 6. Business analysis} Model performance needs to be evaluated against business goals and analyzed to generate business insights. These insights can then be used to eliminate unproductive projects or scope out new projects. Because this step is closely related to the first step, it will also be discussed in chapter 9.

    

\end{itemize}


\textbf{Iterative Process}


Developing an ML system is an iterative and, in most cases, never ending process. You do reach the point where you have to put the system into production, but then that system will constantly need to be monitored and updated.

Before deploying my first ML system, I thought the process would be linear and straightforward. I thought all I had to do was to collect data, train a model, deploy that model, and be done. However, I soon realized that the process looks more like a cycle with a lot of back and forth between different steps.

For example, here is one workflow that you might encounter when building an ML model to predict whether an ad should be shown when users enter a search query.



\begin{itemize}
    \item Choose a metric to optimize. For example, you might want to optimize for impressions -- the number of times an ad is shown.
\end{itemize}


Collect data and obtain labels.
Engineer features.
Train models.
During error analysis, you realize that errors are caused by wrong labels, so you relabel data.
Train model again.
During error analysis, you realize that your model always predicts that an ad shouldn’t be shown, and the reason is because 99.99% of the data you have is no-show (an ad shouldn’t be shown for most queries). So you have to collect more data of ads that should be shown.
Train model again.
Model performs well on your existing test data, which is by now two months ago. But it performs poorly on the test data from yesterday. Your model has degraded, so you need to collect more recent data.
Train model again.
Deploy model.
Model seems to be performing well but then the business people come knocking on your door asking why the revenue is decreasing. It turns out the ads are being shown but few people click on them. So you want to change your model to optimize for clickthrough rate instead.
Go to step 1.










\section{Interview Questions}

1. Explain supervised, unsupervised, weakly supervised, semi-supervised, and active learning.

2. Occam's razor states that when the simple explanation and complex explanation both work equally well, the simple explanation is usually correct. How do we apply this principle in ML?

3. What’s the risk in empirical risk minimization?, Why is it empirical? How do we minimize that risk?


\section{Recommended Readings}

\begin{itemize}
    \item Hands-on Machine Learning with Scikit-learn, keras and tensorflow (chapter 1) 

    \item The Master Algorithm by Peter Domingos 
    
    \item Deep Learning by Ian Goodfellow (chapter 5) 
    
    \item Few useful things to know about machine learning by Peter Domingos 

    \item \href{https://alexandruburlacu.github.io/posts/2022-07-05-neptuneai-automl}{Automl}

    \item \href{https://openreview.net/pdf?id=BZ5a1r-kVsf}{Fun read - A Path Towards Autonomous Machine Intelligence by Yann LeCun}

    

    
    %\item 
\end{itemize}




\subsection{ML Startups - General Readings}

- \href{http://www.shivonzilis.com/ }{The Current State of Machine Intelligence 3.0}

- \href{https://techcrunch.com/2018/09/06/measuring-ai-startups-by-the-right-yardstick/  }{Measuring AI Startups by the right yardstick}

- \href{https://alexandruburlacu.github.io/posts/2022-07-05-neptuneai-automl}{AutoML Solutions: What I Like and Don’t Like About AutoML as a Data Scientist
}
 
- \href{https://a16z.com/2017/03/17/machine-learning-startups-data-saas/}{Competing with big companies:}

%- Microsoft Machine Learning Acceletor: 

%http://www.intellify.com.au/, https://picnet.com.au/artificial-intelligence/, https://altis.com.au/,  https://altis.com.au, https://analytics8.com.au/case-studies/


\section{Code (Python Implementation of ISLR)}

- \url{https://github.com/tdpetrou/Machine-Learning-Books-With-Python/tree/master/Introduction\%20to\%20Statistical\%20Learning}

- \url{https://github.com/JWarmenhoven/ISLR-python}



\chapter{ML Fundamentals}


\section{Linear Regression}



Parametric methods estimate $f$$ by assuming its linear and estimate a set of parameters. 

The goal in linear regression is to build a system that can take a vector $\boldsymbol{x} \in \mathbb{R}^{n}$ as input and predict the value of a scalar $y \in \mathbb{R}$ as its output. The output of linear regression is a linear function of the input. 




\subsection{Example: Linear regression with multiple variables}

To make our housing example more interesting, lets consider a slightly richer dataset in which we also know the number of bedrooms in each house: \\

\begin{tabular}{c|c|c} 
Living area $\left(\right.$ feet $\left.^{2}\right)$ & #bedrooms & Price $(1000 \$ \mathrm{~s})$ \\
\hline 2104 & 3 & 400 \\
1600 & 3 & 330 \\
2400 & 3 & 369 \\
1416 & 2 & 232 \\
3000 & 4 & 540 \\
$\vdots$ & $\vdots$ & $\vdots$
\end{tabular}

Here, the $x^{\prime}$ s are two-dimensional vectors in $\mathbb{R}^{2}$. For instance, $x_{1}^{(i)}$ is the living area of the $i$-th house in the training set, and $x_{2}^{(i)}$ is its number of bedrooms. (In general, when designing a learning problem, it will be up to you to decide what features to choose, so if you are out in Portland gathering housing data, you might also decide to include other features such as whether each house has a fireplace, the number of bathrooms, and so on. We'll say more about feature selection later, but for now lets take the features as given.)
\textbf{To perform supervised learning, we must decide how we're going to represent functions/hypotheses $h$ in a computer. }As an initial choice, lets say we decide to approximate $y$ as a linear function of $x$ :
$$
h_{\theta}(x)=\theta_{0}+\theta_{1} x_{1}+\theta_{2} x_{2}
$$

Here, the $\theta_{i}$ 's are the parameters (also called weights) parameterizing the space of linear functions mapping from $\mathcal{X}$ to $\mathcal{Y}$. When there is no risk of confusion, we will drop the $\theta$ subscript in $h_{\theta}(x)$, and write it more simply as $h(x)$. To simplify our notation, we also introduce the convention of letting $x_{0}=1$ (this is the intercept term), so that

$$
h(x)=\sum_{i=0}^{n} \theta_{i} x_{i}=\theta^{T} x
$$

where on the right-hand side above we are viewing $\theta$ and $x$ both as vectors, and here $n$ is the number of input variables (not counting $x_{0}$ ). we can add an intercept term $b$, often called the bias parameter of the affine transformation 

Now, given a training set, how do we pick, or learn, the parameters $\theta$ ? One reasonable method seems to be to make $h(x)$ close to $y$, at least for the training examples we have.
 
To formalize this, we will define a function that measures, for each value of the $\theta$ 's, how close the $h\left(x^{(i)}\right)$ 's are to the corresponding $y^{(i)}$ 's. We define the \textbf{cost function}:
$$
J(\theta)=\frac{1}{2} \sum_{i=1}^{m}\left(h_{\theta}\left(x^{(i)}\right)-y^{(i)}\right)^{2}
$$

- $h_{\theta}$ means model has been parameterized by the vector $\theta$

Next our goal is to minimize this cost function
$$
\underset{\theta_{0}, \theta_{1}}{\operatorname{minimize}} J\left(\theta_{0}, \theta_{1}\right)
$$


\noindent
\textbf{Terminology}

\begin{itemize}
    \item Objective function is the most general term for any function that you optimize during training. MLE is a type of objective function (which you maximize)

    \item Loss function is usually a function defined on a data point, prediction and label, and measures the penalty. 

    \item Cost function is usually more general. It might be a sum of loss functions over your training set plus some model complexity penalty (regularization). For example One way of measuring the performance of the model is to compute the mean squared error of the model on the test set. If $\hat{\boldsymbol{y}}^{(\text {test }) }
$ gives the predictions of the  model on the test set, then the mean squared error is given by

$$
\mathrm{MSE}_{\text {test }}=\frac{1}{m} \sum_{i}\left(\hat{\boldsymbol{y}}^{(\text {test })}-\boldsymbol{y}^{(\text {test })}\right)_{i}^{2}
$$
Intuitively, one can see that this error measure decreases to 0 when $\hat{\boldsymbol{y}}^{(\text {test })}=\boldsymbol{y}^{\text {(test) }}$. We can also see that
$$
\mathrm{MSE}_{\mathrm{test}}=\frac{1}{m}\left\|\hat{\boldsymbol{y}}^{(\mathrm{test})}-\boldsymbol{y}^{(\mathrm{test})}\right\|_{2}^{2}
$$
so the error increases whenever the Euclidean distance between the predictions and the targets increases.



\end{itemize}




\section{Training the model}

\begin{itemize}
  \item Training a model means searching for a combination of model parameters that minimize a cost function. It is a serach in the model's parameter space: the more parameters a model has, the more dimensions this space has and the harder the search is: searching for a needle in a 300-dimensional haystack is a much trickeir than in 3 dimensions. 

  \item To make a machine learning algorithm, we need to design an algorithm that will improve the weights $\boldsymbol{w}$ in a way that reduces MSE $_{\text {test }}$ when the algorithm is allowed to gain experience by observing a training set $\left(\boldsymbol{X}^{(\text {train })}, \boldsymbol{y}^{(\text {train })}\right) .$ 
  
  \item One intuitive way of doing this is just to minimize the mean squared error on the training set, MSE $_{\text {train }} .$
  
  \item There are two ways to train this particular regression model. One is using a direct 'closed form' equation that directly computes the model parameters that best fit the model to the training set. The other method uses an iternative optimization approach called Gradient Descent that gradually tweaks the model parameters to minimize the cost function over the training set.   
\end{itemize}

\subsection{Method 1: The normal Equation}

 To minimize MSE $_{\text {train }}$, we can simply solve for where its gradient is $\mathbf{0}$ : $$ \begin{array}{c}
\nabla_{\boldsymbol{w}} \mathrm{MSE}_{\text {train }}=0 \\
\Rightarrow \nabla_{\boldsymbol{w}} \frac{1}{m}\left\|\hat{\boldsymbol{y}}^{(\text {train })}-\boldsymbol{y}^{(\text {train })}\right\|_{2}^{2}=0
\end{array}
$$




\begin{equation}
\begin{array}{c}
\Rightarrow \frac{1}{m} \nabla_{\boldsymbol{w}}\left\|\boldsymbol{X}^{(\operatorname{train})} \boldsymbol{w}-\boldsymbol{y}^{(\text {train })}\right\|_{2}^{2}=0 \\
\Rightarrow \nabla_{\boldsymbol{w}}\left(\boldsymbol{X}^{(\text {train })} \boldsymbol{w}-\boldsymbol{y}^{(\text {train })}\right)^{\top}\left(\boldsymbol{X}^{(\operatorname{train})} \boldsymbol{w}-\boldsymbol{y}^{(\text {train })}\right)=0 \\
\Rightarrow \nabla_{\boldsymbol{w}}\left(\boldsymbol{w}^{\top} \boldsymbol{X}^{(\operatorname{train}) \top} \boldsymbol{X}^{(\operatorname{train})} \boldsymbol{w}-2 \boldsymbol{w}^{\top} \boldsymbol{X}^{(\operatorname{train}) \top} \boldsymbol{y}^{(\text {train })}+\boldsymbol{y}^{(\operatorname{train}) \top} \boldsymbol{y}^{(\text {train })}\right)=0 \\
\Rightarrow 2 \boldsymbol{X}^{(\operatorname{train}) \top} \boldsymbol{X}^{(\operatorname{train})} \boldsymbol{w}-2 \boldsymbol{X}^{(\operatorname{train}) \top} \boldsymbol{y}^{(\operatorname{train})}=0 \\
\Rightarrow \boldsymbol{w}=\left(\boldsymbol{X}^{(\operatorname{train}) \top} \boldsymbol{X}^{(\operatorname{train})}\right)^{-1} \boldsymbol{X}^{(\operatorname{train}) \top} \boldsymbol{y}^{(\operatorname{train})}
\end{array}
\end{equation}

\subsection{Method 2: Gradient Descent}




\textbf{Introduction}

Gradient Descent is a generic optimization algorithm capable of finding optimal solutions to a wide range of problems.  Gradient descent is an iterative optimization algorithm for finding the minimum of a function. By reading modern literature on machine learning, you often encounter references to gradient descent or stochastic gradient descent. These are two most frequently used optimization algorithms used in cases where the optimization criterion is differentiable.

\textbf{How it works intuitively?}

Suppose you are lost in the mountains in a dense fog, and you can only feel the slope of the ground below your feet. A good strategy to get to the bottom of the valley quickly is to go donhill in the direction of the steepest slope. This is exactly what Gradient Descent does: it measures the \textbf{local gradient of the error function with regard to the parameter vector (theta) }and it goes in the direction of descending gradient. Once the gradient is zero, you have reached a minimum! 

Conceretely, you start by filling (theta) with random values (also known as random initlization). Then you improve it gradually, taking one baby step at a time, each step attempting to decrease the cost function (MSE) until the algorithm converges to a minimum. 

Learning rate hyperparameter: If the learning rate is too small, then the algorithm will have to go through many iterations to converge, which will take a long time. On the other hand, if the learning rate it too high, you might jump across the valley and end up on the other side, possible even higher up than you were before. Also note that not all cost functions look nice, regular bowls. There may be holes, ridges, plateus, and all sorts of irregular terrains, making convergence to the minimum difficult. Fortunately, the MSE cost function for linear regression is a covex function. So this means there are no local minima, just one gloabl minimum. If the features have different scales, it becomes a elongated scale. r many models, such as logistic regression or SVM, the optimization criterion is convex. Convex functions have only one minimum, which is global. Optimization criteria for neural networks are not convex, but in practice even finding a local minimum local minimum suffices.

\textbf{Technical explanation:}
 
To find a local minimum of a function using gradient descent, one starts at some random point and takes steps proportional to the negative of the gradient (or approximate gradient) of the function at the current point. Gradient descent can be used to find optimal parameters for linear and logistic regression, SVM and also neural networks which we consider later. 



\textbf{Stochastic gradient descent}

Stochastic gradient descent (SGD) is a version of the algorithm that speeds up the computation by approximating the gradient using smaller batches (subsets) of the training data. SGD itself has various “upgrades”. Adagrad is a version of SGD that scales – for each parameter according to the history of gradients. As a result, – is reduced for very large gradients and vice-versa. Momentum is a method that helps accelerate SGD by orienting the gradient descent in the relevant direction and reducing oscillations. In neural network training, variants of SGD such as RMSprop and Adam, are most frequently used


\section{Classification Problem}

As an example, let's pretend we work at a zoo where we have a spreadsheet that contains information about the traits of different animals. We want to use discriminative learning in order to categorize the species of a given animal.

\begin{tabular}{|l|l|l|}
\hline Features & Indepent Variable (X) & Informative columns like num\_legs, has\_wings, has\_shell. \\
\hline Label & Dependent Variable $(y)$ & The species column that we want to predict. \\
\hline
\end{tabular}



\begin{itemize}
    \item In this type of task, the computer program is asked to specify which of $k$ categories some input belongs to. To solve this task, the learning algorithm is usually asked to produce a function $f: \mathbb{R}^{n} \rightarrow\{1, \ldots, k\} .$ When $y=f(\boldsymbol{x})$, the model assigns an input described by vector $\boldsymbol{x}$ to a category identified by numeric code $y$. There are other variants of the classification task, for example, where $f$ outputs a probability distribution over classes.
 
    \item Training data $\mathcal{D}_{n}$ is in the form of a set of pairs $\left\{\left(x^{(1)}, y^{(1)}\right), \ldots,\left(x^{(n)}, y^{(n)}\right)\right\}$ where $x^{(i)}$ represents an object to be classified, most typically a d-dimensional vector of real and/or discrete values, and $y^{(i)}$ is an element of a discrete set of values. The $y$ values are sometimes called target values.

    \item A classification problem is binary or two-class if $y^{(i)}$ is drawn from a set of two possible values; otherwise, it is called multi-class. i.e Binary classification refers to predicting one of two classes and multi-class classification involves predicting one of more than two classes. Multi-label classification involves predicting one or more classes for each example and imbalanced classification refers to classification tasks where the distribution of examples across the classes is not equal.


    \item The goal in a classification problem is ultimately, given a new input value $\chi^{(n+1)}$, to predict the value of $y^{(n+1)}$

    \item Classification problems are a kind of supervised learning, because the desired output (or class) $\mathrm{y}^{(\mathfrak{i})}$ is specified for each of the training examples $\chi^{(\mathfrak{i})}$.
    

\end{itemize}





\textbf{Classification with missing inputs:} Classification becomes more challenging if the computer program is not guaranteed that every measurement inits input vector will always be provided. To solve the Classification task, the learning algorithm only has to define a single function mapping from a vector input to a categorical output. When some of the inputs may be missing,rather than providing a single Classification function, the learning algorithm must learn a set of functions.



\section{Logistic Regression}

With linear regression, we fit the line using "least squares". In other words, we find the line that minimizes the sum of the squares of these residuals. 

Logistic regression predicts whether something is True or False, instead of predicting something continuous like size.

 instead of fitting a line to the data, logistic regression fits an "S" shaped "logistic function". The curve goes from 0 to $1 \ldots$ ...and that means that the curve tells you the probability that a mouse is obese based on its weight.

In other words, just like linear regression, logistic regression can work with continuous data (like weight and age) and discrete data (like genotype and astrological sign).

Logistic regression's ability to provide probabilities and classify new samples using continuous and discrete measurements makes it a popular machine learning method.

In summary... Logistic regression can be used to classify
samples .and it can use different types of data (like size and/or genotype) to do that classification...


\section{Windom from the trenches}

Few things to consider when taking on a Machine learning project: 

\begin{enumerate}
    \item State the assumptions:Create a list of assumptions about the problem and it’s phrasing. These may be rules of thumb and domain specific information that you think will get you to a viable solution faster. It can be useful to highlight questions that can be tested against real data because breakthroughs and innovation occur when assumptions and best practice are demonstrated to be wrong in the face of real data. 

    \item It can also be useful to highlight areas of the problem specification that may need to be challenged, relaxed or tightened.For example: The specific words used in the tweet matter to the model. The specific user that retweets does not matter to the model. The number of retweets may matter to the model. Older tweets are less predictive than more recent tweets.

    \item  What Type of data Do they have?- Identify ML problem category. Do they have enough data for models to generalize?

    \item What accuracy would they be happy with?- How will you measure the project/model’s success?   - What would the acceptance test look like?

    \item Model Deployment(Serving) or Presenting Results:How will the model be deployed? On Cloud, on Premises?

\end{enumerate}


\section{Interview Questions:}

\textbf{What makes a classification problem different from a regression problem? Can a classification problem be turned into a regression problem and vice versa?}

\textbf{Regression:} In this type of task, the computer program is asked to predict a numerical value given some input. To solve this task, the learning algorithm is asked to output a function \begin{equation} f: \mathbb{R}^{n} \rightarrow \mathbb{R} \end{equation} This type of task is similar to classification, except that the format of output is different. \\

\begin{equation}
\text { Regression is like classification, except that } y^{(\mathfrak{i})} \in \mathbb{R}^{k} \text { . }
\end{equation}
    

\chapter{EDA, Visualization and Sampling}

\section{Data Types}


The four basic types of data are:

\begin{itemize}
    \item Nominal data: categorical data with no inherent ordering between the categories. For example, a “pet type” variable could consist of the classes {dog, cat, rabbit}, and there is no relative ordering between these two types, they are just different discrete values.

    \item Ordinal data: categorical data with an inherent ordering, but where the “differences” between categories has no strictly numerical meaning. The canonical example here are survey responses with responses such as: {strong disagree, slightly disagree, neutral, slightly agree, strongly agree}. The important character here is that although there is a clear ordering between these types, there is no sense in which the difference between slightly agree and strongly agree is the “same” as the difference between neutral and slightly agree.

    \item Interval data: numeric data, that is, data that can be mapped to a “number line”; the important aspect in contrast with ordinal data, though, is not the “discrete versus continuous differentiation (integer values can be considered interval data, for instance), but the fact that relative differences in interval data have meaning. A classical example is temperature (in Fahrenheit or Celsius, a point which we will emphasize more shortly): here the differences between temperatures have a meaning: 10 and 15 degrees are separated by the same amount as 15 and 20 (this property is so inherent to numerical data that it almost seems strange to emphasize it). On the other hand, interval data encompasses instances where the zero point has “no real meaning”; what this means in practice is that the ratio between two data points has no meaning. Twenty degrees Farenheit is not “twice as hot” in any meaningful sense than 10 degrees; and certainly not infinitely hotter than zero degrees.

    \item Ratio data: also numeric data, but where the ratio between measurements does have some meaning. The classical example here is temperature Kelvin. Obviously just like temperature Fahrenheit or Celsius, this is describing the basic phenomenon of temperature, but unlike the previous cases, zero Kelvin has a meaning in terms of molecular energy in a substance (i.e., that there is none). This means that ratios have a real meaning: a substance at 20 degrees Kelvin has twice as much kinetic energy at the molecular level as that substance as 10 degrees Kelvin. 


\end{itemize}

 At a course level, the first two data types can simply be viewed as “categorial” data (data taking on discrete values), whereas the later are “numerical” (taking real-valued numbers, though with the caveat that some level of discretization is acceptable even in numeric data, as long as the notion of differences are properly preserved. Indeed, must of the later discussion on visualization will fall exactly along the categorial/real differentiation, though we will highlight some of the considerations when taking in account nominal vs. categorial and interval vs. ratio.



This can also be viewed in terms of the “allowable operations” we can perform on two data point of the different types:

\begin{itemize}
    \item Nominal data: $=, \neq$. All we can do with nominal data is compare two data points and test if they are equal or not.
    \item Ordinal data: $=, \neq,<,>$. In addition to checking for equality, we can also compare the ordering of different data points.
    \item  Interval data: $=, \neq,<,>,-$. We have all the operations of ordinal data, but can also compute exact numeric differences between two data points. Depending on context we may say that addition is allowed too, but sometimes addition actually assumes a zero point and so would only apply to ratio data.
    \item Ratio data: $=, \neq,<, 0,-,+, \div$ We have all the operations of interval data, but addition is now virtually always allowed, and we can additional perform division (to determine ratios) between different data points.

\end{itemize}


\section{Exploratory Data Analysis - Introduction}

\textit{The greatest value of a picture is when it forces us to notice what we never expected to see. -John Tukey}

In the bigger picture Data-driven science, we start by collecting a data set of reasonable size, and then looking for patterns that ideally will play the role of hypotheses for future analysis. Exploratory data analysis is the search for patterns and trends in a given data set. The goal of exploratory data analysis is to get you thinking about your data and reasoning about your question. i.e to make sure we have the right data, any problems with the dataset, determining if what we answer our desired question and get rough idea of what the answer will look like.

People are not very good at looking at a column of numbers or a whole spreadsheet and then determining important characteristics of the data. They find looking at numbers to be tedious, boring, and/or overwhelming. Exploratory data analysis techniques have been devised as an aid in this situation. Most of these techniques work in part by hiding certain aspects of the data while making other aspects more clear.
Overall, The goal of exploratory analysis is to examine or explore the data and find relationships that weren’t previously known. Exploratory analyses explore how different measures might be related to each other but do not confirm that relationship as causitive. After the basic exploratory analysis we can pause and think that if our question needs refinement or if we need to collect more or new data.

Exploratory data analysis is about looking carefully at your dataset and identify any errors in data collection processing, finding violations of statistical assumptions, and suggesting interesting hypotheses.
As first step we could answer the following questions: Who constructed this dataset, when, and why? what is size of the data and the description of the various columns?

EDA always precedes formal (confirmatory) data analysis. EDA is useful for:

\begin{itemize}

    \item Detection of mistakes
    \item Checking of assumptions
    \item Determining relationships among the explanatory variables
    \item Assessing the direction and rough size of relationships between explanatory and outcome variables
    \item Preliminary selection of appropriate models of the relationship between an outcome variable and one or more explanatory variables.

\end{itemize}



\section{Summary of EDA - Common Steps}


\subsection{Before EDA:}

\begin{itemize}
    \item Check the size and type of data
    \item See if the data is in appropriate formart - Convert the data to a format you can easily manupulate (without changing the data itself)
    \item Sample a test set, set it aside and never look at it
\end{itemize}



\subsection{EDA Methods:}

\begin{itemize}
    \item EDA method is either non-graphical or graphical.
    
    \item Each method is either univariate or multivariate (usually just bivariate). Overall,the four types of EDA are univariate non-graphical, multivariate nongraphical, univariate graphical, and multivariate graphical.
    
    \item Non-graphical methods generally involve calculation of summary statistics, while graphical methods obviously summarize the data in a diagrammatic or pictorial way.

    \item Univariate methods look at one variable (data column) at a time, while multivariate methods look at two or more variables at a time to explore relationships. Usually our multivariate EDA will be bivariate (looking at exactly two variables), but occasionally it will involve three or more variables. It is almost always a good idea to perform univariate EDA on each of the components of a multivariate EDA before performing the multivariate EDA.

\end{itemize}


\subsection{EDA Common Steps}

\begin{itemize}
    \item Grab a copy of the data and Read in your data using appropriate python or R libraries 
    \item Check the packaging e.g check the number of rows and columns and see they match the dataset description.
    \item Study Each Attribute and its characteristics (Name, Type, \% of missing values, noisy, usefulness, type of distribution)
    \item Look at the top and the bottom of your data e.g using head() and tail() commands to get a sense of what the rows look like
    \item \textbf{Perform Summary statistics.} e.g Tukey’s five number summary is a great start for numerical values, consisting of the extreme values (max and min), plus the median and quartile elements
    \item \textbf{Visualize the data:} Perform Pairwise correlations, Plots of distributions to identify correlations and identify
    classes.
    \item Formulate your question e.g Are air pollution levels higher on the east coast than on the west coast? For Supervised Machine Learning, identify the target attribute
    \item Study how you would solve the problem manually
    \item Identify the promising transformations you may want to apply and make plans to collect more of different data (if needed and if possible) all of this is an an iterative process to get at the truth and answer our questions.

\end{itemize}
    

\subsection{Recommended Readings:}

\begin{itemize}
    \item Andrew Abela. Advanced Presentations by Design: Creating Communication that Drives Action. Pfeiffer, 2nd edition, 2013.
    \item The Art of Data Science, A Guide for Anyone Who Works with Data by Roger D. Peng and Elizabeth Matsui
    \item Chapter 4, Experimental Design and Analysis by Howard J. Seltman
    \item Section 1, Hands-On Exploratory Data Analysis with Python
    \item Chapter 1, Practical Statistics for Data Scientists, 2nd Edition


\end{itemize}


%\section{Fundamental Statistics Concepts}




%\%textbf{Statistical basic concepts:}
%(Expectation and Mean, Variance and Standard deviation, Covarriance and Correlation, Median Quartile, Interquartile range, Percentile/quantile,Mode )

%- DS field guide book
%- Practical data science book 
%https://www.youtube.com/@statquest/playlists

%knowing which methods are suitable for which type of data



\section{Non-Graphical Methods for EDA}

\subsection{Central tendency} The central tendency or “location” of a distribution has to do with typical or middle values. The common, useful measures of central tendency are the statistics called (arithmetic) mean, median, and sometimes mode. 


%![](https://paper-attachments.dropbox.com/s_478598AFA2F5777FB9289D2A6B80C2413B0ADE29B5C1C858EBCBF98AFD0D611D_1611725655490_Screen+Shot+2021-01-27+at+4.33.58+pm.png)


\begin{enumerate}
    \item \textbf{Mean} For any symmetrically shaped distribution (i.e., one with a symmetric histogram or pdf or pmf) the mean is the point around which the symmetry holds. For non-symmetric distributions, the mean is the “balance point”: if the histogram is cut out of some homogeneous stiff material such as cardboard, it will balance on a fulcrum placed at the mean. 
    \item \textbf{Mode:} The median is another measure of central tendency. The sample median is the middle value after all of the values are put in an ordered list. If there are an even number of values, take the average of the two middle values. (If there are ties at the middle, some special adjustments are made by the statistical software we will use. In unusual situations for discrete random variables, there may not be a unique median.). 
    
    \item \textbf{Mode:} A rarely used measure of central tendency is the mode, which is the most likely or frequently occurring value. More commonly we simply use the term “mode” when describing whether a distribution has a single peak (unimodal) or two or more peaks (bimodal or multi-modal). In symmetric, unimodal distributions, the mode equals both the mean and the median. In unimodal, skewed distributions the mode is on the other side of the median from the mean. In multi-modal distributions there is either no unique highest mode, or the highest mode may well be unrepresentative of the central tendency.

    \item \textbf{Weighted average:} A weighted average is a type of MEAN, where some values in the data set are given more influence than others. Each value to be averaged is given a weight, representing the importance of that value in the average. Weighted averages are important when you are dealing with frequencies or distributions, or when working with data that’s unequal in some way.


\end{enumerate}

\subsubsection{Some useful Tips:}

\textbf{Mean vs Median: } For unimodal skewed (asymmetric) distributions, the mean is farther in the direction of the “pulled out tail” of the distribution than the median is. Therefore, for many cases of skewed distributions, the median is preferred as a measure of central tendency. For example, according to the US Census Bureau 2004 Economic Survey, the median income of US families, which represents the income above and below which half of families fall, was \$43,318. This seems a better measure of central tendency than the mean of \$60,828, which indicates how much each family would have if we all shared equally. \\


\textbf{Robustness:} The median has a very special property called robustness. A sample statistic is “robust” if moving some data tends not to change the value of the statistic. The median is highly robust, because you can move nearly all of the upper half and/or lower half of the data values any distance away from the median without changing the median. More practically, a few very high values or very low values usually have no effect on the median.


The most common measure of central tendency is the mean. For skewed distribution or when there is concern about outliers, the median may be preferred.

Tips:

- MEAN Does your data have a continuous distribution that’s relatively symmetrical? Use MEAN (often referred to as just ‘average’).

- MEDIAN Does your data contain significant outliers? Use MEDIAN - it's less influenced by this.

- MODE represents the most common value in a dataset. If you’re dealing with Nominal data (non-numeric categories like “industry vertical”), MODE is the only appropriate average to use

\subsection{Spread and Variability} 


Several statistics are commonly used as a measure of the spread of a distribution, including variance, standard deviation, and interquartile range. Spread is an indicator of how far away from the center we are still likely to find data values.

The variance and standard deviation are two useful measures of spread. The variance is the mean of the squares of the individual deviations. The standard deviation is the square root of the variance. For Normally distributed data, approximately 95\% of the values lie within 2 sd of the mean.

\textbf{Skewness and kurtosis:} Skewness is a measure of asymmetry. Kurtosis is a more subtle measure of peakedness compared to a Gaussian distribution.



\begin{itemize}
    \item  \textbf{Mode} The most commonly occuring score
   
    \item  \textbf{Median} The score corresponding to the point having 50\% OF the observations below it when the observations are arranged in a numerical order. 
    
    \item \textbf{Mean} The sum of the scores divided by the number of scores
   
    \item \textbf{Dispersion:} The degree to which individual data points are distributed around the mean. 
    
    \item \textbf{Range:} The distance from the lowest to the highest score
    
    \item \textbf{Variance}  it measures the degree of dispersion of data around the sample's mean. It is calculated by taking the differences between each number in the data set and the mean, then squaring the differences to make them positive, and finally dividing the sum of the squares by the number of values in the data set.  Sum of the squared deviations about the mean, divided by N-1. 
    
    \item \textbf{Standard deviation} Although variance its is important, its hard to talk about the variance intuitively. Hence we take The square root of the variance. known as the SD. A low standard deviation indicates that the values tend to be close to the mean (also called the expected value) of the set, while a high standard deviation indicates that the values are spread out over a wider range. 
    
    \item \textbf{Expected value} The long range average of a statistic over repeated samples

    \item \textbf{Covarriance and Correlation}

    \item \textbf{Median Quartile}

    \item \textbf{Interquartile range} The range of the middle 50\% of the observations. 

    \item \textbf{Percentile/Quantile,Mode}


\end{itemize}





\section{Data Visualization }

knowing which methods are suitable for which type of data

Visualisation is a cornerstone of data science and one of the most powerful tools for Exploratory data analysis. Therefore it is important to understand the key principles behind data visualisation before applying your favourite statistical method. Mastering visualisation is a hallmark of an effective data scientist to understand the relevant features of the data at hand.

Edward Tufte has given a lot of thought to this subject and has written several book on this topic. However, our focus in this post is on “visualisation for data exploration” which is different from “visualisation for presentation” and doesn't concerns itself with conveying information/knowledge to a broader audience.

%Data can be more than just storytelling. Some charts can convince decision makers in the effectiveness of certain interventions and others can increase public awareness. e.g This charts shows clear effectiveness of vaccines.



\textbf{Principles of Plotting:}

How to decide which visualization tool is right for you?

\begin{itemize}
    \item Scatter plots — numeric x numeric: If both dimensions of the data are numeric, the most natural first type of plot to consider is the scatter plot: plotting points that simply correspond to the different coordinates of the data.

    \item Line plots - numeric x numeric (sequential): For two dimensional data where one of the dimensions in naturally sequential (this comes up often, for instance, when monitoring a time series, so that one of the two dimensions of each data point is time).

    \item Box and whiskers and violin plots - categorical x numeric: When one dimension is numeric and one is categorical

    \item Heat map and bubble plots — categorical x categorical: When both dimensions of our 2D data are categorical, we have even less information to use.

\end{itemize}

\subsection{Correlation}

It is common to compute correlation coefficients to measure the strength of the relationship between two variables. The correlation between variables $x$ and $y$ is given by
$$
r=\frac{\sum\left(x_t-\bar{x}\right)\left(y_t-\bar{y}\right)}{\sqrt{\sum\left(x_t-\bar{x}\right)^2} \sqrt{\sum\left(y_t-\bar{y}\right)^2}} .
$$
The value of $r$ always lies between $-1$ and 1 with negative values indicating a negative relationship and positive values indicating a positive relationship. 

\subsection{Histograms}

For 1D numerical data (either interval or ratio), the histogram is the natural generalization of a bar chart. Histograms show frequency counts as well, but do so by lumping together ranges in the numerical data (the ranges are typically chosen automatically, i.e., to linearly span over the entire range of the data set). We can construct a histogram using the plt.hist() command as in the following code.

Essentially, what is being done here is that the entire range of the data is divided into bins number of equal-sized bins, spanning the entire range of the data. We then count the number of points that fall into each bin, and generate a bar chart over this binned data. There is some skill to picking the number of bins versus the number of samples in the data, but this can be done by simply trying out a few different values and seeing which produces the most understandable figure.

Histograms are the workhorse of exploratory data analysis. They are indispensable for quickly understanding the distribution of values that a particular feature takes on. And because the distances between points on the x axis do indeed make sense for interval or ratio data, it also is not incorrect to create a line plot showing the overall shape of the distribution.


\subsection{Box and whiskers Plot}

A graphical representation of the dispersion of a sample. 



\subsection{Summary}


Data visualization is arguably the most important tool for exploratory data analysis because the information conveyed by graphical display can be very quickly absorbed and because it is generally easy to recognize patterns in a graphical display


\includegraphics[width=\textwidth]{visualization_techniques.png}






----

[4] Beautiful Evidence by Edward Tufte: http://atc.berkeley.edu/201/readings/Tufte_BE_2006.pdf




\section{Further readings and references}

- \href{}{Fundamental Statistics for Behavioural Sciences (first 6 chapter)}


- \href{http://www.datasciencecourse.org/notes/visualization/}{CMU notes on Data Visualization}

- \href{https://clauswilke.com/dataviz/index.html}{Fundamentals of Data Visualization}

- \href{https://socviz.co/index.html#preface}{Data Visualization - A practical introduction}


\section{Interview Questions}


\textbf{If you have 6 shirts and 4 pairs of pants, how many ways are there to choose 2 shirts and 1 pair of pants?}

\textbf{What is the difference between sampling with vs. without replacement? Name an example of when you would use one rather than the other?}

\textbf{Explain Markov chain Monte Carlo sampling.}

\textbf{ If you need to sample from high-dimensional data, which sampling method would you choose?}

\textbf{Suppose we have a classification task with many classes. An example is when you have to predict the next word in a sentence -- the next word can be one of many, many possible words. If we have to calculate the probabilities for all classes, it’ll be prohibitively expensive. Instead, we can calculate the probabilities for a small set of candidate classes. This method is called candidate sampling. Name and explain some of the candidate sampling algorithms.}
https://www.tensorflow.org/extras/candidate_sampling.pdf






\chapter{Preparing Training Data}

%see: https://docs.google.com/document/d/1eiNkA-e3h5NxYMVGmJLTrkrPArc8r1JTV0GKUOgnZfQ/edit#heading=h.emu98l2wyj1w


\section{Intro: Setting up development and
test sets}


Our predictive algorithm will need samples to learn from as well as samples for evaluating its performance.

So we split our dataset into subsets for these purposes. It's important that the distribution of each subset is representative of the broader population because we want our algorithm to be able to \textbf{generalize}. \\

\begin{tabular}{|l|l|l|}
\hline Train & $67 \%$ & What the algorithm is trained on/ learns from. \\
\hline Validation & $21 \%$ & What the model is evaluated against during training. \\
\hline Test & $12 \%$ & Blind holdout for evaluating the model at the end of training. \\
\hline
\end{tabular}

- Choose dev and test sets from a distribution that reflects what data you expect to get in the future and want to do well on. This may not be the same as your training data's distribution.

- Choose dev and test sets from the same distribution if possible.

- Choose a single-number evaluation metric for your team to optimize. If there are multiple goals that you care about, consider combining them into a single formula (such as averaging multiple error metrics) or defining satisficing and optimizing metrics.

- Machine learning is a highly iterative process: You may try many dozens of ideas before finding one that you're satisfied with.

- Having dev/test sets and a single-number evaluation metric helps you quickly evaluate algorithms, and therefore iterate faster.

- When starting out on a brand new application, try to establish dev/test sets and a metric quickly, say in less than a week. It might be okay to take longer on mature applications.

- The old heuristic of a $70 \% / 30 \%$ train/test split does not apply for problems where you have a lot of data; the dev and test sets can be much less than $30 \%$ of the data.

- Your dev set should be large enough to detect meaningful changes in the accuracy of your algorithm, but not necessarily much larger. Your test set should be big enough to give you a confident estimate of the final performance of your system.

- If your dev set and metric are no longer pointing your team in the right direction, quickly change them: (i) If you had overfit the dev set, get more dev set data. (ii) If the actual distribution you care about is different from the dev/test set distribution, get new $\mathrm{dev} /$ test set data. (iii) If your metric is no longer measuring what is most important to you, change the metric.





\section{Sampling}
For implementing ML based solutions, we face a number of challenges in practice. For example, first we face challenges related to sampling noise and sampling bias. Having non-representative training data for ML training algorithms will lead a scenario where the trained model is unlikely to make accurate predictions.  Data needs vary based on the type of AI project. e.g the requirements for a project building classification intelligence will be different from one providing recommendations or ranking.  We have to consider selecting data that is representative of cases we want to generalize to and one that contains the relevant features needed to learn the desired task. We also need to make sure that we have sufficient data for the problem in question. There are exceptions in scenarios where we can rely on transfer learning. After ensuring the right datasets are available, we have to consider the data quality. To summarize, Data availability broadly translates to having the right quality, quantity, and features. Right quality refers to the fact that the data samples are an accurate reflection of the phenomenon we are trying to model and meet properties such as independent and identically distributed.

It’s often useful to view a data set as a sample from some
distribution or population. How many samples are necessary and sufficient before you can make accurate inferences about the population? How can you estimate what you
don’t know? We’ll also study the Markov Chain Monte Carlo method, by which you can estimate what you cannot compute.

\subsection{Few Useful Concepts:}

\begin{itemize}
    \item \textbf{Population:} Entire collection of events we are interested in e.g lengths of the tails of the cows in the country. Ages as which every girl first began to walk
    
    \item \textbf{Sample:} Set of actual observations: subset of a population
    
    \item \textbf{Statistics:} Numerical Values summarizing sample data
    
    \item \textbf{Parameters} Numerical values summarizing population data
    
    \item \textbf{Random Sample:} A sample in which each member of the population has an equal chance of inclusion. 

    \item \textbf{Bias:} Bias refers to the difference between the measure and some “true” value. A difference between an individual measurement and the true value is called an “error” (which implies the practical impossibility of perfect precision, rather than the making of mistakes). The bias is the average difference over many measurements. Ideally the bias of a measurement process should be zero. 

    \item \textbf{Outliers:} There is no generally recognized formal definition for outlier, but roughly it means values that are outside of the areas of a distribution that would commonly occur. This can also be thought of as sample data values which correspond to areas of the population pdf (or pmf) with low density (or probability). Another common definition of “outlier” consider any point more than a fixed number of standard deviations from the mean to be an “outlier”, but these and other definitions are arbitrary and vary from situation to situation.

\end{itemize}






\subsection{Non-Probability Sampling}
\subsection{Simple Random Sampling}
\subsection{Stratified Sampling	}
\subsection{Weighted Sampling}
\subsection{Importance Sampling	}
\subsection{Reservoir Sampling}




TODO: add 
add tchniqus: \url{https://www.cloudresearch.com/resources/guides/sampling/pros-cons-of-different-sampling-methods/}






\section{Data Challenges in the real-world}
Useful AI systems can be built with anywhere from 100 data points to 1M data points (``so called big data''). But having more data is always a good idea. In mature organizations, AI systems are often time using very sophisticated, multi-year strategies to acquire data, and specific data acquisition strategies are industry and situation-specific. For example, Google and Baidu both have numerous free products that do not monetize but allow them to acquire data that can be monetized elsewhere. But reality is that Training data is often really scarce for many valuable 'ML tasks'. In the real world, 75\% of the challenge in Machine Learning is building the right dataset. AI is great with large scale computer resources, digital data, and efficient algorithms. Without these factors, AI initiatives can fail.  Building high quality models where we're in a position to make accurate predictions in production  going requires a fair amount of effort. For a specialized labeling let's say, for example look at images from radiology and identify different forms of cancer, or a patient's prognosis suddenly one it's much harder to get specialists to do this.


Prof. Peter Norvig, Research and Engineering Director at Google has recommended the following blog post which lays down the concept of the “Machine Learning Surprise” as follows: 

\textit{“Since ML algorithms and optimization are talked about more in literature and media, it is common for people to assume that they play larger roles than they do in the actual implementation process. … Optimizing an ML algorithm takes much less relative effort, but collecting data, building infrastructure, and integration each take much more work. The differences between expectations and reality are profound”. }

Machine learning (ML) models, long the darlings of researchers and practitioners are no longer the center of AI. In fact, models were becoming commodities. It is the training data that would drive progress towards more performant ML models and systems. This dependence of machine learning on labeled training sets is nothing new, and arguably has been the primary driver of new advances for many years (spacemachine.net 2016). But deep learning models are massively more complex than most traditional models–many standard deep learning models today have hundreds of millions of free parameters–and thus require commensurately more labeled training data. These hand-labeled training sets are expensive and time-consuming to create–taking months or years for large benchmark sets, or when domain expertise is required–and cannot be practically repurposed for new objectives. In practice, the cost and inflexibility of hand-labeling such training sets is the key bottleneck to actually deploying machine learning.

\section{Ideas for collecting Additional data}


TODO: update this section form: 

\url{https://docs.google.com/document/d/1eiNkA-e3h5NxYMVGmJLTrkrPArc8r1JTV0GKUOgnZfQ/edit#heading=h.57pb9yyqtl4v}


Understanding ML systems will be helpful in designing and developing them. In this section, we’ll start with the question: how important is data for building intelligent systems?

Progress in the last decade shows that the success of an ML system depends largely on the data it was trained on. Instead of focusing on improving ML algorithms, most companies focus on managing and improving their data.

Despite the success of models using massive amounts of data, many are skeptical of the emphasis on data as the way forward. In the last three years, at every academic conference I attended, there were always some debates among famous academics on the power of mind vs. data. Mind might be disguised as inductive biases or intelligent architectural designs. Data might be grouped together with computation since more data tends to require more computation.

In theory, you can both pursue intelligent design and leverage large data and computation, but spending time on one often takes time away from another.

On the mind over data camp, there’s Dr. Judea Pearl, a Turing Award winner best known for his work on causal inference and Bayesian networks. The introduction to his book, “The book of why”, is entitled “Mind over data,” in which he emphasizes: “Data is profoundly dumb.” In one of his more controversial posts on Twitter, he expressed his strong opinion against ML approaches that rely heavily on data and warned that data-centric ML people might be out of job in 3-5 years.

[QUOTE]
“ML will not be the same in 3-5 years, and ML folks who continue to follow the current data-centric paradigm will find themselves outdated, if not jobless. Take note.”
[/QUOTE]

There’s also a milder opinion from Professor Christopher Manning, Director of the Stanford Artificial Intelligence Laboratory, who argued that huge computation and a massive amount of data with a simple learning algorithm create incredibly bad learners. The structure allows us to design systems that can learn more from fewer data.

Many people in ML today are on the data over mind camp. Professor Richard Sutton, a professor of computing science at the University of Alberta and a distinguished research scientist at DeepMind, wrote a great blog post in which he claimed that researchers who chose to pursue intelligent designs over methods that leverage computation will eventually learn a bitter lesson.

[QUOTE]
“The biggest lesson that can be read from 70 years of AI research is that general methods that leverage computation are ultimately the most effective, and by a large margin. … Seeking an improvement that makes a difference in the shorter term, researchers seek to leverage their human knowledge of the domain, but the only thing that matters in the long run is the leveraging of computation.”
[/QUOTE]

When asked how Google search was doing so well, Peter Norvig, Google’s Director of Search, emphasized the importance of having a large amount of data over intelligent algorithms in their success: “We don’t have better algorithms. We just have more data.”

Dr. Monica Rogati, Former VP of Data at Jawbone, argued that data lies at the foundation of data science, as shown in Figure 1-3. If you want to use data science, a discipline of which machine learning is a part of, to improve your products or processes, you need to start with building out your data, both in terms of quality and quantity. Without data, there’s no data science.


Figure 1-3: The data science hierarchy of needs (Monica Rogati, 2017)

The debate isn’t about whether finite data is necessary, but whether it’s sufficient. The term finite here is important, because if we had infinite data, we can just look up the answer. Having a lot of data is different from having infinite data.

Regardless of which camp will prove to be right eventually, no one can deny that data is essential, for now. Both the research and industry trends in the recent decades show the success of machine learning relies more and more on the quality and quantity of data. Models are getting bigger and using more data. Back in 2013, people were getting excited when the One Billion Words Benchmark for Language Modeling was released, which contains 0.8 billion tokens. Six years later, OpenAI’s GPT-2 used a dataset of 10 billion tokens. And another year later, GPT-3 used 500 billion tokens.












\subsection{Data-Centric AI}

There's a lot of excitement around understanding how to put machine learning to work on real use-cases. Data-Centric AI embodies a particular point of view around how this progress can happen: by focusing on making it easier for practitioners to understand, program and iterate on datasets, instead of spending time on models.

Models are becoming a commodity was heretical for a few reasons. First, people often think of data as a static thing. After all, data literally means “that which is given”. For most ML people, they download an off-the-shelf dataset, drop it into a PyTorch dataloader, and plug-and-play: losses go down, accuracy goes up, and the data is a mere accessory. But to an engineer in the wild, the training data is never “that which is given”. It is the result of a process — usually a dirty, messy process that is critical and underappreciated. For successful applications In applications, teams often take time to clean and merge data. We have to engineer it. 

\textit{Despite the excitement, it is clear to us that success or failure to a level usable in applications we cared about—in medicine, at large technology companies or even pushing the limits on benchmarks—wasn’t really tied to models per se. That is, the advances were impressive, but they were hitting diminishing returns. You can see this in benchmarks, where most of the progress after 2017 is fueled by new advances in augmentations, weak supervision, and other issues of how you feed machines data. In round numbers, ten points of accuracy were due to those—while (by and large) model improvements were squeaking out a few tenths in accuracy points.
}
Overall, AI is driven by data—not code and Training data is the new new oil. There is a data-centric research agenda inside AI. It’s intellectually deep, and it has been lurking at the core of AI progress for a while. Perhaps by calling it out we can make even more progress on an important viewpoint.


\subsection{Labeling}


TODO: https://mitpress.mit.edu/9780262047074/machine-learning-from-weak-supervision/

Many traditional lines of research in machine learning are similarly motivated by the insatiable appetite of modern machine learning models for labeled training data. We start by drawing the core distinction between these other approaches and weak supervision at a high-level: weak supervision is about leveraging higher-level and/or noisier input from subject matter experts (SMEs).

For simplicity, we can start by considering the categorical classification setting: we have data points \(x \in X\) that each have some label \(y \in Y = {1,...,K}\), and we wish to learn a classifier \(h : X \rightarrow Y\). For example, \(X\) might be mammogram images, and \(Y\) a tumor grade classification. We choose a model class \(h_\theta\)–for example a CNN–and then pose our learning problem as one of parameter estimation, with our goal being to minimize the expected loss, e.g. on a new unseen test set of labeled data points \((x_i, y_i)\). In the standard supervised learning setting, we then would get a training set of data points with ground-truth labels \(T = {(x_1, y_1), ..., (x_N, y_N)}\)–traditionally, hand-labeled by SME annotators, e.g. radiologists in our example–define a loss function \(L(h(x), y)\), and minimize the aggregate loss on the training set using an optimization procedure like SGD.

The problem is that this is expensive: for example, unlike grad students, radiologists don’t generally accept payment in burritos and free T-shirts! Thus, many well-studied lines of work in machine learning are motivated by the bottleneck of getting labeled training data:



\textbf{Vertical farming case study: Tackling Supervised Classification Problems}


Labelled data is essential for applying any supervised learning technique. Labels can indicate if a plant has a condition of interest. However, Labelling plant data is not straightforward and has a lot of subtleties that one has to be wary of. Expensive stress region annotations by domain experts are needed to implement image-based phenotyping solutions.  Building a stress image dataset containing diverse species is time-consuming and involves expert involvement(e.g plant pathologists). The ideal dataset will contain images taken under varying and challenging conditions. Diversity can be measured in terms of images of all types: illumination, complex background, resolution, size, and orientation. 

Generally speaking, there are two approaches to create labelled datasets:

- Manually labelling:  This is the most obvious straight forward way where a human observer would carefully go through the plant data and label the plant for the condition of interest.

- Writing an algorithm that could automatically label plants. Here Rule-based systems are used as a manual method where a deterministic approach for automatically labelling data is adopted. Such algorithms rely on a list a set of rules, created on the basis of expert knowledge to label a plant for the condition of interest. Examples of rule-based algorithms for various 'human' diseases can be found on PheKB. However, we note that maintaining a large rule-base is challenging as solutions based on that will have high predictive value but poor recall. We also note that recently, machine learning based labelling techniques are being used to solve this problem. Here, we can imagine setting up a supervised learning problem where we are able to automatically learn rules for labelling a plants using only few manually labelled patients.


\textbf{In active learning}, the goal is to make use of subject matter experts more efficiently by having them label data points which are estimated to be most valuable to the model. Traditionally, applied to the standard supervised learning setting, this means selecting new data points to be labeled–for example, we might select mammograms that lie close to the current model decision boundary, and ask radiologists to label only these. However, we could also just ask for weaker supervision pertinent to these data points, in which case active learning is perfectly complementary with weak supervision; as one example of this, see (Druck, Settles, and McCallum 2009).

In the semi-supervised learning setting, we have a small labeled training set and a much larger unlabeled data set. At a high level, we then use assumptions about smoothness, low dimensional structure, or distance metrics to leverage the unlabeled data (either as part of a generative model, as a regularizer for a discriminative model, or to learn a compact data representation); for a good survey see (Chapelle, Scholkopf, and Zien 2009). More recent methods use adversarial generative (Salimans et al. 2016), heuristic transformation models (Laine and Aila 2016), and other generative approaches to effectively help regularize decision boundaries. Broadly, rather than soliciting more input from subject matter experts, the idea in semi-supervised learning is to leverage domain- and task-agnostic assumptions to exploit the unlabeled data that is often cheaply available in large quantities.

In the standard transfer learning setting, our goal is to take one or more models already trained on a different dataset and apply them to our dataset and task; for a good overview see (Pan and Yang 2010). For example, we might have a large training set for tumors in another part of the body, and classifiers trained on this set, and wish to apply these somehow to our mammography task. A common transfer learning approach in the deep learning community today is to “pre-train” a model on one large dataset, and then “fine-tune” it on the task of interest. Another related line of work is multi-task learning, where several tasks are learned jointly (Caruna 1993; Augenstein, Vlachos, and Maynard 2015). Some transfer learning approaches take one or more pre-trained models (potentially with some heuristic conditioning of when they are each applied) and use these to train a new model for the task of interest; in this case, we can actually consider transfer learning as a type of weak supervision.
The above paradigms potentially allow us to avoid asking our SME collaborators for additional training labels. But what if–either in addition, or instead–we could ask them for various types of higher-level, or otherwise less precise, forms of supervision, which would be faster and easier to provide? For example, what if our radiologists could spend an afternoon specifying a set of heuristics or other resources, that–if handled properly–could effectively replace thousands of training labels? This is the key practical motivation for weak supervision approaches, which we describe next.

\textbf{
Weak Supervision:} A Simple Definition
In the weak supervision setting, our objective is the same as in the supervised setting, however instead of a ground-truth labeled training set we have:

Unlabeled data \(X_u = x_1, …, x_N\);
One or more weak supervision sources \(\tilde{p}_i(y \vert x), i=1:M\) provided by a human subject matter expert (SME), such that each one has:
A coverage set \(C_i\), which is the set of points \(x\) over which it is defined
An accuracy, defined as the expected probability of the true label \(y^*\) over its coverage set, which we assume is \(\lt 1.0\)
In general, we are motivated by the setting where these weak label distributions serve as a way for human supervision to be provided more cheaply and efficiently: either by providing higher-level, less precise supervision (e.g. heuristic rules, expected label distributions), cheaper, lower-quality supervision (e.g. crowdsourcing), or taking opportunistic advantage of existing resources (e.g. knowledge bases, pre-trained models). These weak label distributions could thus take one of many well-explored forms:

\textbf{Weak Labels:} The weak label distributions could be deterministic functions–in other words, we might just have a set of noisy labels for each data point in \(C_i\). These could come from crowd workers, be the output of heuristic rules \(f_i(x)\), or the result of distant supervision (Mintz et al. 2009), where an external knowledge base is heuristically mapped onto \(X_u\). These could also be the output of other classifiers which only yield MAP estimates, or which are combined with heuristic rules to output discrete labels.
\textbf{Constraints:} We can also consider constraints represented as weak label distributions. Though straying outside of the simple categorical setting we are considering here, the structured prediction setting leads to a wide range of very interesting constraint types, such as physics-based constraints on output trajectories (Stewart and Ermon 2017) or output constraints on execution of logical forms (Clarke et al. 2010; Guu et al. 2017), which encode various forms of domain expertise and/or cheaper supervision from e.g. lay annotators.

Distributions: We might also have direct access to a probability distribution. For example, we could have the posterior distributions of one or more weak (i.e. low accuracy/coverage) or biased classifiers, such as classifiers trained on different data distributions as in the transfer learning setting. We could also have one or more user-provided label or feature expectations or measurements (Mann and McCallum 2010; Liang, Jordan, and Klein 2009), i.e. an expected distribution \(p_i(y)\) or \(p_i(y\vert f(x))\) (where \(f(x)\) is some feature of \(x\)) provided by a domain expert as in e.g. (Druck, Settles, and McCallum 2009).
Invariances: Finally, given a small set of labeled data, we can express functional invariances as weak label distributions–e.g., extend the coverage of the labeled distribution to all transformations of \(t(x)\) or \(x\), and set \(p_i(y\vert t(x)) = p_i(y\vert x)\). In this way we view techniques such as data augmentation as a form of weak supervision as well.

\section{Data Augmentation}



Data augmentation is a family of techniques that are used to increase the amount of training data. Traditionally, these techniques are used for tasks that have limited training data, such as in medical imaging projects. However, in the last few years, they have shown to be useful even when we have a lot of data because augmented data can make our models more robust to noise and even adversarial attacks.

Data augmentation has become a standard step in many computer vision tasks and is finding its way into natural language processing (NLP) tasks. The techniques depend heavily on the data format, as image manipulation is different from text manipulation. In this section, we will cover three main types of data augmentation: simple label-preserving transformations, perturbation, which is a term for “adding noises”, and data synthesis. In each type, we’ll go over examples for both computer vision and NLP.

\textbf{Simple Label-Preserving Transformations}

In computer vision, the simplest data augmentation technique is to randomly modify an image while preserving its label. You can modify the image by cropping, flipping, rotating, inverting (horizontally or vertically), erasing part of the image, and more. This makes sense because a rotated image of a dog is still a dog. Common ML frameworks like PyTorch and Keras both have support for image augmentation. According to Krizhevsky et al., in their legendary AlexNet paper, “the transformed images are generated in Python code on the CPU while the GPU is training on the previous batch of images. So these data augmentation schemes are, in effect, computationally free.”

In NLP, you can randomly replace a word with a similar word, assuming that this replacement wouldn’t change the meaning or the sentiment of the sentence. Similar words can be found either with a dictionary of synonymous words, or by finding words whose embeddings are close to each other in a word embedding space.

Original sentences

I’m so happy to see you.
Generated sentences
I’m so glad to see you.
I’m so happy to see y’all.
I’m very happy to see you.

Table 3-8: Three sentences generated from an original sentence by replacing a word 
with another word with similar meaning.

This type of data augmentation is a quick way to double, even triple your training data.
Perturbation
Perturbation is also a label-preserving operation, but because sometimes, it’s used to trick models into making wrong predictions, I thought it deserves its own section.

Neural networks, in general, are sensitive to noise. In the case of computer vision, this means that by adding a small amount of noise to an image can cause a neural network to misclassify it.  Su et al. showed that 67.97% of the natural images in Kaggle CIFAR-10 test dataset and 16.04% of the ImageNet test images can be misclassified by changing just one pixel (See Figure 3-11).

Using deceptive data to trick a neural network into making wrong predictions is called adversarial attacks. Adding noise to samples to create adversarial samples is a common technique for adversarial attacks. The success of adversarial attacks is especially exaggerated as the resolution of images increases.

Adding noisy samples to our training data can help our models recognize the weak spots in their learned decision boundary and improve their performance,. Noisy samples can be created by either adding random noise or by a search strategy. Moosavi-Dezfooli et al. proposed an algorithm, called DeepFool, that finds the minimum possible noise injection needed to cause a misclassification with high confidence. This type of augmentation is called adversarial augmentation.

Adversarial augmentation is less common in NLP (an image of a bear with randomly added pixels still looks like a bear, but adding random characters to a random sentence will render it gibberish), but perturbation has been used to make models more robust. One of the most notable examples is BERT, where the model chooses 15% of all tokens in each sequence at random, and chooses to replace 10% of the chosen tokens with random words. For example, given the sentence “my dog is hairy” and the model randomly replaces “hairy” with “apple”, the sentence becomes “my dog is apple”. So 1.5% of all tokens might result in nonsensical meaning. Their ablation studies show that a small fraction of random replacement gives their model a small performance boost.

In chapter 5, we’ll go over how to use perturbation not just as a way to improve your model’s performance, but also a way to evaluate its performance.


reference: \url{https://docs.google.com/document/d/1N7dRx5zwnyXCl0H-VuKpAHTbMBiZXRHozQ7pScSdz1s/edit}



%\textbf{\textit{Adaptive Learning:}} In a scenario where large amounts of unlabeled data are available but the task of labelling is hard or infeasible, adaptive learning methods are very useful. Adaptive Learning methods adaptively select the most informative samples for labeling for the highest improvement in test accuracy. The goal of Adaptive Learning is to achieve maximum predictive performance under a fixed labeling budget, which makes it useful for applications with limited data availability. Many Adaptive Learning methods have been proposed in the literature with different heuristics to reduce the amount of labeling needed for training machine learning (ML) models for classification tasks. 

%To get started with adaptive learning we select a small sample chosen randomly initially for labeling.  In the second phase, this dataset is then used to train a machine  learning model. In the last phase, a random batch of data from the remaining unlabeled data is adaptively selected using an acquisition function for labelling by human domain experts. The acquisition function helps in selecting the most useful samples in the unlabeled dataset for improving model performance. Overall, the retraining of the model continues until one of two termination criteria is met – a desired performance threshold of the model is achieved, or the labeling budget is exhausted \cite{nagasubramanian2020useful}.

%Weak supervision, domain adaptation, domain randomization, synthetic dataset creation, and transfer learning are some of the other promising methods that we plan on experimenting with, in order to reduce the amount of labeling needed. Combining super-resolution with conventional images can also be used used to enhance the spatial resolution of desired images \cite{nagasubramanian2020useful}.  Lastly, the \textit{Amazon SageMaker Ground Truth} provides a service for managing the labelling, including two features. One is annotation consolidation, which combines different people’s annotation task results into one high-fidelity label. The second one is automated data labeling, which utilizes machine learning to label portions of the provided data automatically.

%For applying deep learning techniques, appropriate architecture needs to be selected based on the size and type of available data and the relative complexity of the problem. Training isolated models specific to individual plant species will prove to be a tedious process. Transfer learning models  are usually trained for a specific task, such as a computer vision in plant phenotyping problem, and can subsequently be used for training similar models for related problems/tasks. Transfer learning is useful where the learnt model (and associated representations) from one problem is transferred to a somewhat dissimilar but relevant problem.  The type of transfer learning to use on new dataset depends on two factors: the size of the new dataset and its similarity with the the data used to train the original model. \cite{singh2018deep}. In the context of deep learning, transfer learning may involve (i) reusing the entire model (architecture + parameters) for the new task that can involve a new dataset, (ii) only reusing the model architecture/hyperparameters, but learning the parameters from scratch with the new dataset, and (iii) reusing a model but after fine tuning the model parameters based on the new dataset – this process is often referred to as \textit{domain adaptation}.


%Overall, a dedicated effort is required in the long term where we build specific datasets and ensure the availability of multiple types, scales, and resolution of data. This entails annotation, augmentation, outlier rejection, standardization, and normalization of data. We note that similarity of training data to original data leads to easier transferability and in practice it is advised that we initialize weights from a pre-trained model, as it has the benefit of allowing quicker/robust fine-tuning throughout the entire network. 






\section{Interview Questions}

https://huyenchip.com/ml-interviews-book/contents/7.2-sampling-and-creating-training-data.html


\section{Useful tools}

we want to cover Tools & methodologies for accelerating open-source dataset iteration:}


Example use-cases of such tools are as follows: 

- Tools that quantify and accelerate time to source and prepare high quality data

- Tools that ensure that the data is labeled consistently, such as label consensus

- Tools that make improving data quality more systematic

- Tools that automate the creation of high quality supervised learning training data from low quality resources, such as forced alignment in speech recognition

- Tools that produce consistent and low noise data samples, or remove labeling noise or inconsistencies from existing data

- Tools for controlling what goes into the dataset and for making high level edits efficiently to very large datasets, e.g. adding new words, languages, or accents to speech datasets with thousands of hours

- Search methods for finding suitably licensed datasets based on public resources

- Tools for creating training datasets for small data problems, or for rare classes in the long tail of big data problems

- Tools for timely incorporation of feedback from production systems into datasets

- Tools for understanding dataset coverage of important classes, and editing them to cover newly identified important cases
Dataset importers that allow easy combination and composition of existing datasets

- Dataset exporters that make the data consumable for models and interface with model training and inference systems such as webdataset.

- System architectures and interfaces that enable composition of dataset tools such as, MLCube, Docker, Airflow

\textbf{Algorithms for working with limited labeled data and improving label efficiency:
}
Data selection techniques such as active learning and core-set selection for identifying the most valuable examples to label.
Semi-supervised learning, few-shot learning, and weak supervision methods for maximizing the power of limited labeled data.
Transfer learning and self-supervised learning approaches for developing powerful representations that can be used for many downstream tasks with limited labeled data.
Novelty and drift detection to identify when more data needs to be labeled.



https://datacentricai.org/neurips21/

- Clean lab 
- https://datacentricai.org/



\chapter{Feature Engineering}

In 2014, the paper Practical Lessons from Predicting Clicks on Ads at Facebook claimed that having the right features is the most important thing in developing their ML models. Since then, many of the companies  have discovered time and time again that once they have a workable model, having the right features tends to give them the biggest performance boost compared to clever algorithmic techniques such as hyperparameter tuning. State-of-the-art model architectures can still perform poorly if they don’t use a good set of features.

Due to its importance, a large part of many ML engineering and data science jobs is to come up with new useful features. In this chapter, we will go over common techniques and important considerations with respect to feature engineering. 

We will end the chapter discussing how to engineer good features, taking into account both the feature importance and feature generalization.

\section{Common Feature Engineering Operations}


\subsection{Handling Missing Values}

Missing not at random (MNAR): 
Missing not at random (MNAR): 
Missing completely at random (MCAR):


There are three types of missing values. The official names for these types are a little bit confusing so we’ll go into detailed examples to mitigate the confusion.

\textbf{Missing not at random (MNAR): }when the reason a value is missing is because of the value itself. In this example, we might notice that respondents of gender “B” with higher income tend not to disclose their income. The income values are missing for reasons related to the values themselves.

\textbf{Missing at random (MAR):} when the reason a value is missing is not due to the value itself, but due to another observed variable. In this example, we might notice that age values are often missing for respondents of the gender “A”, which might be because the people of gender A in this survey don’t like disclosing their age.

\textbf{Missing completely at random (MCAR): }when there’s no pattern in when the value is missing. In this example, we might think that the missing values for the column “Job” might be completely random, not because of the job itself and not because of another variable. People just forget to fill in that value sometimes for no particular reason. However, this type of missing is very rare. There are usually reasons why certain values are missing, and you should investigate.


\textbf{Deletion}

When I ask candidates about how to handle missing values during interviews, many tend to prefer deletion, not because it’s a better method, but because it’s easier to do.

One way to delete is column deletion: if a variable has too many missing values, just remove that variable. For example, in the example above, over 50\% of the values for the variable “Marital status” are missing, so you might be tempted to remove this variable from your model. The drawback of this approach is that you might remove important information and reduce the accuracy of your model. Marital status might be highly correlated to buying houses, as married couples are much more likely to be homeowners than single people.

Another way to delete is row deletion: if an example has missing value(s), just remove that example from the data. This method can work when the missing values are completely at random (MCAR) and the number of examples with missing values is small, such as less than 0.1\%. You don’t want to do row deletion if that means 10\% of your data examples are removed.

However, removing rows of data can also remove important information that your model needs to make predictions, especially if the missing values are not at random (MNAR). For example,  you don’t want to remove examples of gender B respondents with missing income because 

\textbf{Imputation}

Even though deletion is tempting because it’s easy to do, deleting data can lead to losing important information or cause your model to be biased. If you don’t want to delete missing values, you will have to impute them, which means “fill them with certain values.” Deciding which “certain values” to use is the hard part.

One common practice is to fill in missing values with their defaults. For example, if the job is missing, you might fill it with an empty string “”. Another common practice is to fill in missing values with the mean, median, or mode (the most common value). For example, if the temperature value is missing for a data example whose month value is July, it’s not a bad idea to fill it with the median temperature of July.

Both practices work well in many cases, but sometimes, they can cause hair-splitting bugs. One time, in one of the projects I was helping with, we discovered that the model was spitting out garbage because the app’s front-end no longer asked users to enter their age, so age values were missing, and the model filled them with 0. But the model never saw the age value of 0 during training, so it couldn’t make reasonable predictions.

In general, you might not want to fill missing values with possible values, such as filling the missing number of children with 0 — 0 is a possible value for the number of children. It makes it hard to distinguish between people for whom you don’t have children information and people who don’t have children.

Multiple techniques might be used at the same time or in sequence to handle missing values for a particular set of data. Regardless of what techniques you use, one thing is certain: there is no perfect way to handle missing values. With deletion, you risk losing important information or accentuating biases. With imputation, you risk adding noise to your data, or worse, data leakage. If you don’t know what data leakage is, don’t panic, we’ll cover it in the Data Leakage section of this chapter.


\subsection{Scaling}

Consider the task of predicting whether someone will buy a house in the next 12 months, and the data is shown in Table 4-2. The values of the variable Age in our data go between 20 and 40, whereas the values of the variable Annual Income go between 10,000 and 150,000. When we input these two variables into an ML model, it won’t understand that 150,000 and 40 represent different things. It will just see them both as numbers, and because the number 150,000 is much bigger than 40, it might give it more importance, regardless of which variable is actually more useful for the predicting task.

During data processing, it’s important to scale your features so that they’re in similar ranges. This process is called feature scaling. This is one of the simplest things you can do that often result in a performance boost for your model. Neglecting to do so can cause your model to make gibberish predictions, especially with classical algorithms like gradient-boosted trees and logistic regression.

An intuitive way to scale your features is to get each feature to be in the range [0, 1]. Given a variable x, its values can be rescaled to be in this range using the following formula.

$$x' = \frac{x - \min(x)}{\max(x) - \min(x)}$$

You can validate that if x is the maximum value, the scaled value x’ will be 1. If x is the minimum value, the scaled value x’ will be 0.

If you want your feature to be in an arbitrary range [a, b] — empirically, I find the range [-1, 1] to work better than the range [0, 1] — you can use the following formula.

$$x' = a + \frac{(x - \min(x))(b - a)}{\max(x) - \min(x)}$$

Scaling to an arbitrary range works well when you don’t want to make any assumptions about your variables. If you think that your variables might follow a normal distribution, it might be helpful to normalize them so that they have zero-mean and unit variance. This process is called standardization.

with x being the mean of variable x, and  being its standard deviation.
$$x' = \frac{x - \bar{x}}{\sigma}$$

with $\bar{x}$ being the mean of variable x, and $\sigma$ being its standard deviation.

In practice, ML models tend to struggle with features that follow a skewed distribution. To help mitigate the skewness, a technique commonly used is log transformation: apply the log function to your feature. An example of how the log transformation can make your data less skewed is shown in Figure 4-3. While this technique can yield performance gain in many cases, it doesn’t work for all cases and you should be wary of the analysis performed on log-transformed data instead of the original data.

There are two important things to note about scaling. One is that it’s a common source of data leakage, (this will be covered in greater detail in the Data Leakage section). Another is that it often requires global statistics — you have to look at the entire or a subset of training data to calculate its min, max, or mean. During inference, you reuse the statistics you had obtained during training to scale new data. If the new data has changed significantly compared to the training, these statistics won’t be very useful. Therefore, it’s important to retrain your model often to account for these changes. We’ll discuss more on how to handle changing data in production in the section on continual learning in Chapter 8.

\subsection{Details on Feature Scaling}


 Scaling changes the bounds of the data, and can be useful, for example, when you are working with data in different units.
 

Next we need to  Make sure features are on a similar scale. Target value scaling is generally not required. This normalization also helps features start on equal footing and prevent gradient explosion.
 Two standard approaches are min-max scaling and standardization

In min-max scaling, values are shifted and rescaled so that they end up ranging from 0 to 1. We subtract the min value and dividing by the max minus the min. We an use sci-kit learn's MinMax scaler for this. 

Standardization involves subtracting the mean value and then dividing by the standard deviation so that the resulting distribution has unit variance. Sci-kit learn provides StandardScalr for standardization. 


\begin{itemize}
    \item Make'sure'features'are'on'a'similar'scale
    \item Get every feature into approximately a $-1 \leq x_{i} \leq 1$ r range
    \item Mean normalization:  
    \item Replace $x_{i}$ with $x_{i}-\mu_{i}$ to make features have approximately zero mean
    \item Redundant'features'(linearly'dependent)'
    \item curse of dimensionality is a myth
    \item It has been proven that if learning rate $\alpha$ is sufficiently small, then $J(\theta)$ will decrease on every iteration. Andrew Ng recommends decreasing $\alpha$ by multiples of $3 .$
    
\end{itemize}

\subsection{Discretization}

Imagine that we’ve built a model with the data in Table 4-2. During training, our model has seen the annual income values of 150000, 50000, 100000, 50000, 60000, and 10000. During inference, our model encounters an example with an annual income of 9000.50. 

Intuitively, we know that  \$9000.50 a year isn’t much different from \$10,000/year, and we want our model to treat both of them the same way. But the model doesn’t know that. Our model only knows that 9000.50 is different from 10000, and will treat them differently.

Discretization is the process of turning a continuous feature into a discrete feature. This process is also known as quantization. This is done by creating buckets for the given values. For annual income, you might want to group them into three buckets as follows.

Lower income: less than \$35,000/year
Middle income: between \$35,000 and \$100,000/year
Upper income: more than \$100,000/year

Now, instead of having to learn an infinite number of possible incomes, our model can focus on learning only three categories, which is a much easier task to learn.

Even though by definition, it’s used for continuous features, the same technique can be used for discrete features too. The age variable is discrete, it might still be useful to group them into buckets such as follows.
Less than 18
Between 18 and 22
Between 22 and 30
Between 30 - 40
Between 40 - 65
Over 65

A question with this technique is how to best choose the boundaries of categories. You can try to plot the histograms of the values and choose the boundaries that make sense. In general, common sense, basic quantiles, and sometimes subject matter expertise can get you a long way.


\subsection{Encoding Categorical Features}

Data needs the be numerically encoded before it can be fed into a neural network.


\begin{tabular}{|l|l|l|}
\hline Binarize & Categorical & 1 means presence, 0 means absence. \\
\hline OneHotEncode(OHE) & Categorical & Expand 1 multi-category col into many binary cols. \\
\hline Ordinal & Categorical & [Bad form] Each category assigned an integer $[0,1,2]$. \\
\hline Scale & Continuous & Shrink the range of values between $-1: 1$ or $0: 1$. \\
\hline
\end{tabular}

We’ve talked about how to turn continuous features into categorical features. In this section, we’ll discuss how to best handle categorical features.

People who haven’t worked with data in production tend to assume that categories are static, which means the categories don’t change over time. This is true for many categories. For example, age brackets and income brackets are unlikely to change and you know exactly how many categories there are in advance. Handling these categories is straightforward. You can just give each category a number and you’re done.

However, in production, categories change. Imagine you’re building a recommendation system to predict what products users might want to buy for Amazon. One of the features you want to use is the product brand. When looking at Amazon’s historical data, you realize that there are a lot of brands. Even back in 2019, there were already over 2 million brands on Amazon!

The number of brands is overwhelming but you think: “I can still handle this.” You encode each brand a number, so now you have 2 million numbers from 0 to 1,999,999 corresponding to 2 million brands. Your model does spectacularly on the historical test set, and you get approval to test it on 1\% of today’s traffic.

In production, your model crashes because it encounters a brand it hasn’t seen before and therefore can’t encode. New brands join Amazon all the time. You create a category “UNKNOWN” with the value of 2,000,000 to catch all the brands your model hasn’t seen during training.

Your model doesn’t crash anymore but your sellers complain that their new brands are not getting any traffic. It’s because your model didn’t see the category UNKNOWN in the train set, so it just doesn’t recommend any product of the UNKNOWN brand. Then you fix this by encoding only the top 99\% most popular brands and encode the bottom 1\% brand as UNKNOWN. This way, at least your model knows how to deal with UNKNOWN brands.

Your model seems to work fine for about 1 hour, then the click rate on recommended products plummets. Over the last hour, 20 new brands joined your site, some of them are new luxury brands, some of them are sketchy knockoff brands, some of them are established brands. However, your model treats them all the same way it treats unpopular brands in the training set.

This isn’t an extreme example that only happens if you work at Amazon on this task. This problem happens quite a lot. For example, if you want to predict whether a comment is spam, you might want to use the account that posted this comment as a feature, and new accounts are being created all the time. The same goes for new product types, new website domains, new restaurants, new companies, new IP addresses, and so on. If you work with any of them, you’ll have to deal with this problem.

Finding a way to solve this problem turns out to be surprisingly difficult. You don’t want to put them into a set of buckets because it can be really hard—how would you even go about putting new user accounts into different groups?

One solution to this problem is the hashing trick, popularized by the package Vowpal Wabbit developed at Microsoft. The gist of this trick is that you use a hash function to generate a hashed value of each category. The hashed value will become the index of that category. Because  We’ve talked about how to turn continuous features into categorical features. In this section, we’ll discuss how to best handle categorical features.

People who haven’t worked with data in production tend to assume that categories are static, which means the categories don’t change over time. This is true for many categories. For example, age brackets and income brackets are unlikely to change and you know exactly how many categories there are in advance. Handling these categories is straightforward. You can just give each category a number and you’re done.

However, in production, categories change. Imagine you’re building a recommendation system to predict what products users might want to buy for Amazon. One of the features you want to use is the product brand. When looking at Amazon’s historical data, you realize that there are a lot of brands. Even back in 2019, there were already over 2 million brands on Amazon!

The number of brands is overwhelming but you think: “I can still handle this.” You encode each brand a number, so now you have 2 million numbers from 0 to 1,999,999 corresponding to 2 million brands. Your model does spectacularly on the historical test set, and you get approval to test it on 1\% of today’s traffic.

In production, your model crashes because it encounters a brand it hasn’t seen before and therefore can’t encode. New brands join Amazon all the time. You create a category “UNKNOWN” with the value of 2,000,000 to catch all the brands your model hasn’t seen during training.

Your model doesn’t crash anymore but your sellers complain that their new brands are not getting any traffic. It’s because your model didn’t see the category UNKNOWN in the train set, so it just doesn’t recommend any product of the UNKNOWN brand. Then you fix this by encoding only the top 99\% most popular brands and encode the bottom 1\% brand as UNKNOWN. This way, at least your model knows how to deal with UNKNOWN brands.

Your model seems to work fine for about 1 hour, then the click rate on recommended products plummets. Over the last hour, 20 new brands joined your site, some of them are new luxury brands, some of them are sketchy knockoff brands, some of them are established brands. However, your model treats them all the same way it treats unpopular brands in the training set.

This isn’t an extreme example that only happens if you work at Amazon on this task. This problem happens quite a lot. For example, if you want to predict whether a comment is spam, you might want to use the account that posted this comment as a feature, and new accounts are being created all the time. The same goes for new product types, new website domains, new restaurants, new companies, new IP addresses, and so on. If you work with any of them, you’ll have to deal with this problem.

Finding a way to solve this problem turns out to be surprisingly difficult. You don’t want to put them into a set of buckets because it can be really hard—how would you even go about putting new user accounts into different groups?

One solution to this problem is the hashing trick, popularized by the package Vowpal Wabbit developed at Microsoft. The gist of this trick is that you use a hash function to generate a hashed value of each category. The hashed value will become the index of that category. Because 

You can choose a hash space large enough to reduce the collision. You can also choose a hash function with properties that you want, such as a locality-sensitive hashing function where similar categories (such as websites with similar names) are hashed into values close to each other.

Because it’s a trick, it’s often considered hacky by academics and excluded from ML curricula. But its wide adoption in the industry is a testimonial to how effective the trick is. It’s essential to Vowpal Wabbit and it’s part of the frameworks scikit-learn, TensorFlow, and gensim. It can be especially useful in continual learning settings where your model learns from incoming examples in production. We’ll cover continual learning in Chapter 8.

\textbf{Feature Crossing}

Feature crossing is the technique to combine two or more features to generate new features. This technique is useful to model the non-linear relationships between features. For example, for the task of predicting whether someone will want to buy a house in the next 12 months, you suspect that there might be a non-linear relationship between marital status and number of children, so you combine them to create a new feature “marriage and children” as in Table 4-3.


Because feature crossing helps model non-linear relationships between variables, it’s essential for models that can’t learn or are bad at learning non-linear relationships, such as linear regression, logistic regression, and tree-based models. It’s less important in neural networks, but can still be useful because explicit feature crossing occasionally helps neural networks learn non-linear relationships faster. DeepFM and xDeepFM are the family of models that have successfully leverage explicit feature interactions for recommendation systems and click-through-rate prediction tasks.

A caveat of feature crossing is that it can make your feature space blow up. Imagine feature A has 100 possible values and feature B has 100 possible features, crossing these two features will result in a feature with 100 x 100 = 10,000 possible values. You will need a lot more data for models to learn all these possible values. Another caveat is that because feature crossing increases the number of features models use, it can make models overfit to the training data.



\section{Feature Encoding for Event Data}
Trace Encoding in Process Mining: a survey and benchmarking

https://arxiv.org/abs/2301.02167


\section{Feature encoding for text/vision}


First introduced to the deep learning community in the paper Attention Is All You Need (Vaswani et al., 2017), positional embedding has become a standard data engineering technique for many applications in both computer vision and natural language processing. We’ll walk through an example to show why positional embedding is necessary and how to do it.

Consider the task of language modeling where you want to predict the next token based on the previous sequence of tokens. In practice, a token can be a word, a character, or a subword, and a sequence length can be up to 512 if not larger. However, for simplicity, let’s use words as our tokens and use the sequence length of 8. Given an arbitrary sequence of 8 words, such as, “Sometimes all I really want to do is”, we want to predict the next word.

If we use a recurrent neural network, it will process words in sequential order, which means the order of words is implicitly inputted. However, if we use a model like a transformer, words are processed in parallel, so words’ positions need to be explicitly inputted so that our model knows which word follows which word (“a dog bites a child” is very different from “a child bites a dog”). We don’t want to input the absolute positions:  0, 1, 2, …, 7 into our model because empirically, neural networks don’t work well with inputs that aren’t unit-variance (that’s why we scale our features, as discussed previously in the section Scaling).

If we rescale the positions to between 0 and 1, so 0, 1, 2, …, 7 become 0, 0.143, 0.286, …, 1, the differences between the two positions will be too small for neural networks to learn to differentiate.

A way to handle position embeddings is to treat it the way we’d treat word embedding. With word embedding, we use an embedding matrix with the vocabulary size as its number of columns, and each column is the embedding for the word at the index of that column. With position embedding, the number of columns is the number of positions. In our case, since we only work with the previous sequence size of 8, the positions go from 0 to 7 (see Figure 4-5). 

The embedding size for positions is usually the same as the embedding size for words so that they can be summed. For example, the embedding for the word “food” at position 0 is the sum of the embedding vector for the word “food” and the embedding vector for position 0. This is the way position embeddings are implemented in HuggingFace’s BERT as of August 2021. Because the embeddings change as the model weights get updated, we say that the position embeddings are learned.

Position embeddings can also be fixed. The embedding for each position is still a vector with S elements (S is the position embedding size), but each element is predefined using a function, usually sine and cosine. In the original Transformer paper, if the element is at an even index, use sine. Else, use cosine. See Figure 4-6.

Fixed positional embedding is a special case of what is known as Fourier features. If positions in positional embeddings are discrete, Fourier features can also be continuous. Consider the task involving representations of 3D objects, such as a teapot. Each position on the surface of the teapot is represented by a 3-dimensional coordinate, which is continuous. When positions are continuous, it’d be very hard to build an embedding matrix with continuous column indices, but fixed position embeddings using sine and cosine functions still work.

This is the generalized format for the embedding vector at coordinate v, also called the Fourier features of coordinate v. Fourier features have been shown to improve models’ performance for tasks that take in coordinates (or positions) as inputs.


$$\gamma(v) = [a_1 \cos(2\pi{b_1}^Tv), a_1\sin(2\pi{b_1}^Tv), …, a_m\cos(2\pi{b_m}^Tv), a_m\sin(2\pi{b_m}^Tv)]^T$$
\textbf{More readings}

In this section, we cover the general heursitics of feature engineering. 

- https://machinelearningmastery.com/discover-feature-engineering-how-to-engineer-features-and-how-to-get-good-at-it/
- https://www.youtube.com/watch?v=ko5WEfLqCh8 
- https://trevorstephens.com/kaggle-titanic-tutorial/r-part-4-feature-engineering/
https://machinelearningmastery.com/feature-selection-with-real-and-categorical-data/ 
https://www.oreilly.com/library/view/feature-engineering-for/9781491953235/
https://machinelearningmastery.com/feature-engineering-and-selection-book-review/ 











\section{Practical useful tools}


\subsection{Data Cleaning}


These tools focus on identify errors in datasets, without taking the downstream model or application into account. These include traditional constraint-based data cleaning methods, as well as those that use machine learning to detect and resolve data errors.

\textbf{HoloClean} functional dependencies, quantitative statistics, external information as a single factor-graph model.

\textbf{Raha} uses a library of error detectors, and treats the output of each as a feature in a holistic detection model. It then uses clustering and active learning to train the holistic model with few labels.

\textbf{Picket:} Self-supervised Data Diagnostics for ML 

\textbf{Pipelines:} self-supervision to learn an error detection model.


\subsection{Library Functions}

- Code to perform these transformations
- create transformation functions that allow us to repeatedly perform these functions 


\section{Interview Questions}




\textbf{How do we ensure that the transformation happen in the correct order? }




Transformation Ppeilines and Custom Transformers 
- A sequence of data processing components is called a data pipeline. 
These components can run asychronysly, where each component pulls in a large amount of data, processesing it, and spits ut the result in another data store. 

- Scikitlearn allows us to establish the pipeline so that transformation occur in the right order




\chapter{Model Evaluation and Debugging Strategies}



\section{Introduction}

The evaluation of the regression task involves computing a distance between your predictions and the values of the ground truth. This difference can be evaluated in different ways, for instance by squaring it in order to punish large errors or by applying a log to it in order to penalize predictions of the wrong scale. 

Classification can be binary, multi-class or multi-label. Mulit-label means class predictions are not exclusive and you can predict multiple class ownership. 

For example in the titanic competition submissions are evaluated based on accuracy, the percentage of surviving passengers that you correctly predict. In housing prices competition, its evaluated based on the average difference between your prediction and the ground truth. In the real world we might use multiple metrics.

If we talk about classification problems, the most common metrics used are:

- Accuracy
- Precision (P)
- Recall (R)
- F1 score (F1)
- Area under the ROC (Receiver Operating Characteristic) curve or simply AUC (AUC)
- Log loss
- Precision at k $(\mathrm{P} @ \mathrm{k})$
- Average precision at $\mathrm{k}(\mathrm{AP} @ \mathrm{k})$
- Mean average precision at k (MAP@k)

When it comes to regression, the most commonly used evaluation metrics are:
- Mean absolute error (MAE)
- Mean squared error (MSE)
- Root mean squared error (RMSE)
- Root mean squared logarithmic error (RMSLE)
- Mean percentage error (MPE)
- Mean absolute percentage error (MAPE)
$-R^2$

\section{Decoupling Objectives}

When solving a problem, you might have multiple objectives in mind. Imagine you're building a system to rank items on users' newsfeed. Your original goal is to maximize users' engagement. You want to achieve this goal through the following three objectives.

1. Filter out spam
2. Filter out NSFW content
3. Rank posts by engagement: how likely users will click on it


However, you quickly learned that optimizing for users' engagement alone can lead to questionable ethical concerns. Because extreme posts tend to get more engagements, your algorithm learned to prioritize extreme content ${ }^{12}$. You want to create a more wholesome newsfeed. So you have a new goal: maximize users' engagement while minimizing the spread of extreme views and misinformation. To obtain this goal, you add two new objectives to your original plan.

\begin{enumerate}
    \item Filter out spam
    \item Filter out NSFW content
    \item Filter out misinformation
    \item Rank posts by quality
    \item Rank posts by engagement: how likely users will click on it
\end{enumerate}


Now, objectives 4 and 5 are in conflict with each other. If a post is very engaging but it's of questionable quality, should that post rank high or low?

Let's take a step back to see what exactly each objective does. To rank posts by quality, you first need to predict posts' quality and you want posts' predicted quality to be as close to their actual quality as possible. Essentially, you want to minimize quality\_loss: the difference between each post's predicted quality and its true quality ${ }^3$.

Similarly, to rank posts by engagement, you first need to predict the number of clicks each post will get. You want to minimize engagement\_loss: the difference between each post's predicted clicks and its actual number of clicks.

One approach is to combine these two losses into one loss and train one model to minimize that loss.
$\text { loss }=\alpha \text { quality\_loss }+\beta \text { engagement\_loss }$

You can randomly test out different values of $\alpha$ and $\beta$ to find the values that work best. If you want to be more systematic about tuning these values, you can check out Pareto optimization, "an area of multiple criteria decision making that is concerned with mathematical optimization problems involving more than one objective function to be optimized simultaneously $y^4$ ".

A problem with this approach is that each time you tune $\alpha$ and $\beta$ - for example, if your users' newsfeed's quality goes up but users' engagement goes down, you might want to decrease $\alpha$ and increase $\beta$ - you'll have to retrain your model.

Another approach is to train two different models, each optimizing one loss. So you have two models:
- quality\_model minimizes quality\_loss and outputs the predicted quality of each post.
- engagement_model minimizes engagement\_loss and outputs the predicted number of clicks of each post.

You can combine the outputs of these two models and rank posts by their combined scores: $\alpha$ quality\_score $+\beta$ engagement\_score
Now you can tweak $\alpha$ and $\beta$ without retraining your models!
In general, when there are multiple objectives, it's a good idea to decouple them first because it makes model development and maintenance easier. First, it's easier to tweak your system without retraining models, as explained above. Second, it's easier for maintenance since different objectives might need different maintenance schedules. Spamming techniques evolve much faster than the way post quality is perceived, so spam filtering systems need updates at a much higher frequency than quality ranking systems.






\subsection{Basic terminology:}

\begin{itemize}
    \item \textbf{Loss function} function that is defined on a single data point and considering the prediction of the model and the ground truth for the data point, computes a penalty. 

    \item \textbf{Cost function} takes into account the whole dataset used for training, computing a sum or average over the loss penalties of its data points. It can compromiuse further constraints, such as L1 or L2. Cost function direct

    \item \textbf{Objective function} is the most general term related to the scope of optimization during machine learning training


    
\end{itemize}



\section{Cross-Validation}

cross-validation is a step in the process of building a machine-learning model which helps us ensure that our models fit the data accurately and also ensures that we do not overfit.




The next important type of cross-validation is stratified k-fold. If you have a skewed dataset for binary classification with $90 \%$ positive samples and only $10 \%$ negative samples, you don't want to use random k-fold cross-validation. Using simple k-fold cross-validation for a dataset like this can result in folds with all negative samples. In these cases, we prefer using stratified k-fold cross-validation. Stratified k-fold cross-validation keeps the ratio of labels in each fold constant. So, in each fold, you will have the same $90 \%$ positive and $10 \%$ negative samples. Thus, whatever metric you choose to evaluate, it will give similar results across all folds.




\section{Wisdom on Picking Cost function for regression}






\section{Evaluating Classifiers}

\textbf{Accuracy:} It is one of the most straightforward metrics used in machine learning. It defines how accurate your model is. For the problem described above, if you build a model that classifies 90 images accurately, your accuracy is 90\% or 0.90. If only 83 images are classified correctly, the accuracy of your model is 83\% or 0.83.

if the dataset is skewed, i.e., the number of samples in one class outnumber the number of samples in other class by a lot. In
these kinds of cases, it is not advisable to use accuracy as an evaluation metric as it is not representative of the data. So, you might get high accuracy, but your model will probably not perform that well when it comes to real-world samples, and you won’t be able to explain to your managers why

assumed that chest x-ray images with pneumothorax are positive class (1) and without pneumothorax are negative class (0). if your model incorrectly (or falsely) predicts positive class, it is a false positive. If your model incorrectly (or falsely) predicts negative class, it is a
false negative. 


\subsection{Confusion Matrix}
\subsection{Precision and Recall}


Precision $=T P /(T P+F P)$ 
Recall $=T P /(T P+F N)$

\subsection{Precision/Recall Tradeoff}

\subsection{The ROC Curve}



\section{Debugging strategy}

- When you start a new project, especially if it is in an area in which you are not an expert, it is hard to correctly guess the most promising directions.

- So don't start off trying to design and build the perfect system. Instead build and train a basic system as quickly as possible-perhaps in a few days. Then use error analysis to help you identify the most promising directions and iteratively improve your algorithm from there.

- Carry out error analysis by manually examining $~ 100$ dev set examples the algorithm misclassifies and counting the major categories of errors. Use this information to prioritize what types of errors to work on fixing.

- Consider splitting the dev set into an Eyeball dev set, which you will manually examine, and a Blackbox dev set, which you will not manually examine. If performance on the Eyeball dev set is much better than the Blackbox dev set, you have overfit the Eyeball dev set and should consider acquiring more data for it.

- The Eyeball dev set should be big enough so that your algorithm misclassifies enough examples for you to analyze. A Blackbox dev set of 1,000-10,000 examples is sufficient for many applications.

- If your dev set is not big enough to split this way, just use the entire dev set as an Eyeball dev set for manual error analysis, model selection, and hyperparameter tuning.


\textbf{Techniques for reducing avoidable bias
}If your learning algorithm suffers from high avoidable bias, you might try the following techniques:

- Increase the model size (such as number of neurons/layers): This technique reduces bias, since it should allow you to fit the training set better. If you find that this increases variance, then use regularization, which will usually eliminate the increase in variance.

- Modify input features based on insights from error analysis: Say your error analysis inspires you to create additional features that help the algorithm eliminate a particular category of errors. (We discuss this further in the next chapter.) These new features could help with both bias and variance. In theory, adding more features could increase the variance; but if you find this to be the case, then use regularization, which will usually eliminate the increase in variance.

- Reduce or eliminate regularization (L2 regularization, L1 regularization, dropout): This will reduce avoidable bias, but increase variance.

- Modify model architecture (such as neural network architecture) so that it is more suitable for your problem: This technique can affect both bias and variance.

One method that is not helpful:
- Add more training data: This technique helps with variance problems, but it usually has no significant effect on bias.


\textbf{Techniques for reducing avoidable bias
}

If your learning algorithm suffers from high avoidable bias, you might try the following techniques:

- Increase the model size (such as number of neurons/layers): This technique reduces bias, since it should allow you to fit the training set better. If you find that this increases variance, then use regularization, which will usually eliminate the increase in variance.

- Modify input features based on insights from error analysis: Say your error analysis inspires you to create additional features that help the algorithm eliminate a particular category of errors. (We discuss this further in the next chapter.) These new features could help with both bias and variance. In theory, adding more features could increase the variance; but if you find this to be the case, then use regularization, which will usually eliminate the increase in variance.

- Reduce or eliminate regularization (L2 regularization, L1 regularization, dropout): This will reduce avoidable bias, but increase variance.

- Modify model architecture (such as neural network architecture) so that it is more suitable for your problem: This technique can affect both bias and variance.

One method that is not helpful:
- Add more training data: This technique helps with variance problems, but it usually has no significant effect on bias.

\href{https://nessie.ilab.sztaki.hu/~kornai/2020/AdvancedMachineLearning/Ng_MachineLearningYearning.pdf
}{Machine Learning Yearning 
}

\section{Model Offline Evaluation}

https://docs.google.com/document/d/1_kxf0xRAushBEepKh57mHubWCkVuS37iAqwmJcYFoJE/edit



\section{Model Monitoring: Probalistic Classification Inference Performance  Analysis}


- https://www.lokad.com/probabilistic-forecasting-definition#:~:text=A%20forecast%20is%20said%20to,outcome%20as%20%E2%80%9Cthe%E2%80%9D%20forecast.
- https://medium.com/nerd-for-tech/probability-and-machine-learning-570815bad29d
- https://machinelearningmastery.com/probability-metrics-for-imbalanced-classification/
- https://www.quora.com/What-is-the-probabilistic-classification-in-machine-learning
- https://www.youtube.com/watch?v=e0Pa5PC_8mQ&t=92s&ab_channel=StanfordOnline
- https://medium.com/knowledge-engineering-seminar/evaluating-probabilistic-classifier-roc-and-pr-g-curves-6647c3379d60

- https://towardsdatascience.com/introduction-to-probabilistic-classification-a-machine-learning-perspective-b4776b469453
- https://stats.stackexchange.com/questions/352518/area-under-the-precision-recall-curve-similar-interpretation-to-auroc
- http://ee104.stanford.edu/lectures.html

More:

https://github.com/evidentlyai/evidently
https://snorkel.ai/#

- https://www.fiddler.ai/model-monitoring
- https://christophergs.com/machine%20learning/2020/03/14/how-to-monitor-machine-learning-models/ 
- https://towardsdatascience.com/monitoring-machine-learning-models-in-production-how-to-track-data-quality-and-integrity-391435c8a299 

https://towardsdatascience.com/evaluate-model-performance-with-cumulative-gains-and-lift-curves-1f3f8f79da01

https://towardsdatascience.com/monitoring-ml-models-in-production-768b6a74ee51




\section{Research - Future of Model Evaluation}

Model evaluation is a crucial part of the model development process in machine learning. The goal of evaluation is to understand the quality of a model, and anticipate if it will perform well in the future.

While evaluation is a classical problem in machine learning, data-centric AI approaches have catalyzed a shift towards fine-grained evaluation: moving beyond standard measures of average performance such as accuracy and F1 scores, to measuring performance on particular populations of interest. This enables a more granular understanding of model performance, and gives users a clearer idea of model capabilities. This shift is complementary to a growing interest in understanding model robustness, since access to fine-grained evaluation permits an enhanced ability to build more robust models.

Approaches to fine-grained evaluation include measuring performance on critical data subsets called slices, invariance or sensitivity to data transformations, and resistance to adversarial perturbations. While most evaluation is user-specified, an important line of work found that models often underperform on hidden strata that are missed by model builders in evaluation, which can have profound consequences on our ability to deploy and use models. This motivates future work in automatically discovering these hidden strata, or more generally, finding all possible failure modes of a model by analyzing datasets and models systematically in conjunction.

Another important facet of fine-grained evaluation is data and model monitoring in order to anticipate, measure and mitigate degradations in performance due to distribution shift. This includes identifying and isolating data points that may be considered outliers, estimating performance on unlabeled data that is streaming to a deployed model, and generating rich summaries of how the data distribution may be shifting over time.


One standard assumption for successfully deploying machine learning models is that test time distributions are similar to those encountered and well-represented during training. In reality however, this assumption rarely holds: seldom do we expect to deploy models in settings that exactly match their training distributions. Training models robust to distribution shifts is then another core challenge to improve machine learning in the wild, which we argue can be addressed under a data-centric paradigm.

Here, we broadly categorize attempts to improve robustness to distribution shifts as those addressing (1) subpopulation shift or hidden stratification, (2) domain shift, and (3) shifts from adversarial perturbations.

Under subpopulation shift, training and test-time distributions differ in how well-represented each subpopulation or “data group” is. If certain subpopulations are underrepresented in the training data, then even if these distributions are encountered during training, standard empirical risk minimization (ERM) and “learning from statistical averages” can result in models that only perform well on the overrepresented subpopulations.

One initial real world problem under subpopulation shift came with training models on datasets that exhibit spurious correlations. If a majority of groups exhibit relations between certain features and the target of interest, but these dependencies do not hold for all data, then models may learn non-robust dependencies relying on these “spurious” correlations. If these groups are known, Group DRO can prevent this by focusing optimization on worst-group error.
Another instantiation comes from hidden stratification, where datapoints belonging to the same labeled classes can actually vary in their feature distributions quite a bit. With GEORGE, we learned that despite not being able to generalize to unseen data under all groups, deep neural networks trained with ERM can actually learn separable representations for different groups that share the same label.
Both Group DRO and GEORGE introduced approaches to handle subpopulation shift under real-world instantiations. These methods have inspired additional work related to upsampling estimated groups (LfF, JTT) and using contrastive learning to learn group-invariant representations (CNC - link coming soon).

Beyond subpopulation shift, robustness also features domain shift and adversarial perturbations. Under domain shift, we model test-time data as coming from a completely different domain from the training data. Under distribution shift with adversarial perturbations, test-time data may exhibit corruptions or imperceptible differences in input feature space that prevent trained ERM models from strongly generalizing to the test-time distributions. 


\section{Key Resources}

\begin{itemize}
    \item Chapter 3 - Hands on ML 
    \item Kaggle book (chapter 
    \item Approaching (Almost) Any Machine Learning Problem

\end{itemize}






\chapter{Theory of Learning}




\chapter{Time-Series Modeling}

Forecasting is a common statistical task in business, where it helps to inform decisions about the scheduling of production, transportation and personnel, and provides a guide to long-term strategic planning. Forecasting is about predicting the future as accurately as possible, given all of the information available, including historical data and knowledge of any future events that might impact the forecasts.

Time series forecasting is an important research area for machine learning (ML), particularly where accurate forecasting is critical, including several industries such as retail, supply chain, energy, finance, etc. For example, in the consumer goods domain, improving the accuracy of demand forecasting by 10-20\% can reduce inventory by 5\% and increase revenue by 2-3\%. Current ML-based forecasting solutions are usually built by experts and require significant manual effort, including model construction, feature engineering and hyper-parameter tuning.

Two currently popular approaches to nonlinear time series prediction problems are statistical approaches using ARIMA and machine learning approaches where we use curve fitting. e.g using Artificial Neural Networks.


\section{Data Cleaning and Pre-Processing}

\begin{itemize}
    \item Missing data
    \item Changing the frequency of a time series (that is, upsampling and downsampling)
    \item Smoothing data
    \item Addressing seasonality in data
    \item Preventing unintentional lookaheads
\end{itemize}



\subsection{Smoothing data}

Use smoothing to eliminate unwanted noise or high variance in data.

\subsection{Detrending data}

Removing a trend from the data lets you focus your analysis on the fluctuations in the data about the trend. While trends can be meaningful, others are due to systematic effects, and some types of analyses yield better insight once you remove them. Removing offsets is another, similar type of preprocessing.

\subsection{Time-Domain and Frequency Domain Transformations}

Alternative bases and the Fourier perspective.This topic continues the theme of howa  shift  in  perspective  can  illuminate  otherwise  undetectable  patterns  in  data.   For example, some data has a temporal component (like audio or time-series data).  Other data has locality (nearby pixels of an image are often similar, same for measurements by sensors).  A naive representation of such data might have one point per moment in time or per point in space.  It can be far more informative to transform the data into a dual" representation, which rephrases the data in terms of patterns that occur across  time  or  across  space.   This  is  the  point  of  the  Fourier  transform  and  other similar-in-spirit transforms.










\subsection{Handling Missing Data}

- \textbf{Imputation}: When we fill in missing data based on observations about the entire data set.

- \textbf{Interpolation}: When we use neighboring data points to estimate the missing value. Interpolation can also be a form of imputation.

- \textbf{Deletion of affected time periods}
When we choose not to use time periods that have missing data at all.  \\


Pandas provides Forward fill, Moving average and Interpolation methods. 

We can also impute data with either a rolling mean or median. Known as a moving average, this is similar to a forward fill in that you are using past values to “predict”. missing future values (imputation can be a form of prediction). With a moving aver‐ age, however, you are using input from multiple recent times in the past. There are many situations where a moving average data imputation is a better fit for the task than a forward fill. For example, if the data is noisy, and you have reason to doubt the value of any individual data point relative to an overall mean, you should use a moving average rather than a forward fill.

Interpolation is a method of determining the values of missing data points based on geometric constraints regarding how we want the overall data to behave. For example, a linear interpolation constrains the missing data to a linear fit consistent with known neighboring points.


\subsection{Resampling}

You should take the temporal resolution of the timestamping you receive with a grain of salt, based on domain knowledge about the behavior you are studying and also based on the details you can determine relating to how the data was collected.

For example, imagine you are looking at daily sales data, but you know that in many cases managers wait until the end of the week to report figures, estimating rough daily numbers rather than recording them each day. The measurement error is likely to be substantial due to recall problems and innate cognitive biases. You might consider changing the resolution of your sales data from daily to weekly to reduce or average out this systematic error. Otherwise, you should build a model that factors in the possibility of biased errors across different days of the week. For example, it may be that managers systematically overestimate their Monday performance when they report the numbers on a Friday.


\subsection{Smoothing}

While outlier detection is a topic in and of itself, if you have reason to believe your data should be smoothed, you can do so with a moving average to eliminate measure‐ ment spikes, errors of measurement, or both. Even if the spikes are accurate, they may not reflect the underlying process and may be more a matter of instrumentation problems; this is why it’s quite common to smooth data.
Smoothing data is strongly related to imputing missing data, and so some of those techniques are relevant here as well. For example, you can smooth data by applying a rolling mean, with or without a lookahead, as that is simply a matter of the point’s position relative to the window used to calculate its smoothed value.

\subsection{Group-by Operations}

With time series data, there are groups just as there are with nontemporal data, so group-by-style thinking is still applicable. For example, in cross-sectional data you might take group means based on age, gender, or neighborhood. In time series analysis, you might find exploratory group-by operations useful, such as computing a monthly average or weekly medians. You can also combine temporal and non tempo‐ ral groups, such as the monthly average calorie intake per gender in a population, or the weekly median amount of sleep per age group in a set of hospital patients. There are many ways that grouping data can allow the temporal axis to interact with other relationships within the data. For example:

\section{Exploratory Data Analysis}

The process is the same as what you have performed on non-time-series data to start with.

- Are any of the columns strongly correlated with one another?
- What is the overall mean of an interesting variable? What is its variance?

To answer these, you can use familiar techniques such as plotting, taking summary statistics, applying histograms, and using targeted scatter plots while we take into account the temporal axis. 

There exist Time Series–Specific Exploratory Methods:

\subsection{Stationary}

What it means for a time series to be stationary and a statistical test for stationarity. The first question you will likely ask about a time series is whether it appears to reflect a system that is “stable” or one that is constantly changing. The level of stability, or stationarity, is important to assess because we need to know how much we should expect the system’s long-term past behavior to reflect its long-term future behavior. 

\subsection{Self-correlation}

Once we have assessed the “stability” (this word is not used in a technical sense here) of a time series, we try to determine whether there are internal dynamics. That is, we are looking for self- correlations, answering the fundamental question of how tightly past data, distant or recent, predicts future data. 


What it means to say that a time series correlates with itself and what such a correlation indicates about the underlying dynamics of the time series.  

\subsection{Spurious correlations}

Finally, when we think we have found certain behavioral dynamics within the system, we need to make sure we are not identifying relation‐ ships based on dynamics that do not in any way imply the causal relationships we wish to discover; hence, we must look for spurious correlations.

What it means for a correlation to be spurious and when you should expect to run into spurious correlations. 




\section{Time series patterns}

\subsection{Trend}

A trend exists when there is a long-term increase or decrease in the data. It does not have to be linear. Sometimes we will refer to a trend as “changing direction”, when it might go from an increasing trend to a decreasing trend. There is a trend in the antidiabetic drug sales data shown in Figure 2.2.



\subsection{Seasonal}
A seasonal pattern occurs when a time series is affected by seasonal factors such as the time of the year or the day of the week. Seasonality is always of a fixed and known frequency. The monthly sales of antidiabetic drugs above shows seasonality which is induced partly by the change in the cost of the drugs at the end of the calendar year.

\subsection{Cyclic}

A cycle occurs when the data exhibit rises and falls that are not of a fixed frequency. These fluctuations are usually due to economic conditions, and are often related to the “business cycle”. The duration of these fluctuations is usually at least 2 years.


\subsection{Trend Calculations}

A moving average of order $m$ can be written as
$$
\hat{T}_t=\frac{1}{m} \sum_{j=-k}^k y_{t+j}
$$
where $m=2 k+1$. That is, the estimate of the trend-cycle at time $t$ is obtained by average values of the time series within $k$ periods of $t$. Observations that are nearby in time are also likely to be close in value. Therefore, the average eliminates some of the randomness in the data, leaving a smooth trend-cycle component. We call this an $m$-MA, meaning a moving average of order $m$.

Moving averages (MAs), moving histograms, and time dependent rates calculate time-dependent statistics from time series while giving more importance to recent than to old samples. The exponential MA (EMA) is computed as: 

$$\quad A_i=a \cdot A_{i-1}+(1-a) \cdot X_i$$

$X_i$ is the size of an observed sample at time $i$ and $A_i$ is the time-dependent average at that time. The smoothing parameter $a$ is often set to $a=0.9$, but there is only little insight in the literature about appropriate configuration.

An exponential moving average, EMA, can be thought of as a filter that puts greater importance on most recent values within the aggregation window. There are two main parameters to an EMA, i.e., the alpha decay rate and the window size. The larger the time window or lower the alpha, the slower moving the EMA will be. In other words, the moving average will lag behind the actual time series by more.

\subsection{Trend Acceleration}

Differenced exponential moving averages can be used to detect the velocity, or trend, of the time series. This approach can be applied a second time, on the differenced time series, to detect the acceleration, or the rate of change of the trend.

\section{Predictions and Forecasting}

The thing we are trying to forecast is unknown (or we would not be forecasting it), and so we can think of it as a random variable. For example, the total sales for next month could take a range of possible values, and until we add up the actual sales at the end of the month, we don’t know what the value will be. So until we know the sales for next month, it is a random quantity.

Because next month is relatively close, we usually have a good idea what the likely sales values could be. On the other hand, if we are forecasting the sales for the same month next year, the possible values it could take are much more variable. In most forecasting situations, the variation associated with the thing we are forecasting will shrink as the event approaches. In other words, the further ahead we forecast, the more uncertain we are.

We can imagine many possible futures, each yielding a different value for the thing we wish to forecast When we obtain a forecast, we are estimating the middle of the range of possible values the random variable could take. Often, a forecast is accompanied by a prediction interval giving a range of values the random variable could take with relatively high probability. For example, a 95\% prediction interval contains a range of values which should include the actual future value with probability 95\%.

We will use the subscript $t$ for time. For example, $y_t$ will denote the observation at time $t$. Suppose we denote all the information we have observed as $\mathcal{I}$ and we want to forecast $y_t$. We then write $y_t \mid \mathcal{I}$ meaning "the random variable $y_t$ given what we know in $\mathcal{I}$ ". The set of values that this random variable could take, along with their relative probabilities, is known as the "probability distribution" of $y_t \mid \mathcal{I}$. In forecasting, we call this the forecast distribution.

When we talk about the "forecast", we usually mean the average value of the forecast distribution, and we put a "hat" over $y$ to show this. Thus, we write the forecast of $y_t$ as $\hat{y}_t$, meaning the average of the possible values that $y_t$ could take given everything we know.
Occasionally, we will use $\hat{y}_t$ to refer to the median (or middle value) of the forecast distribution instead.

It is often useful to specify exactly what information we have used in calculating the forecast. Then we will write, for example, $\hat{y}_{t \mid t-1}$ to mean the forecast of $y_t$ taking account of all previous observations $\left(y_1, \ldots, y_{t-1}\right)$. Similarly, $\hat{y}_{T+h \mid T}$ means the forecast of $y_{T+h}$ taking account of $y_1, \ldots, y_T$ (i.e., an $h$-step forecast taking account of all observations up to time $T$ ).
\subsection{Baseline Methods}


\textbf{Average method}

Here, the forecasts of all future values are equal to the average (or "mean") of the historical data. If we let the historical data be denoted by $y_1, \ldots, y_T$, then we can write the forecasts as
$$
\hat{y}_{T+h \mid T}=\bar{y}=\left(y_1+\cdots+y_T\right) / T .
$$
The notation $\hat{y}_{T+h \mid T}$ is a short-hand for the estimate of $y_{T+h}$ based on the data $y_1, \ldots, y_T$.

\textbf{Naive method}

For naïve forecasts, we simply set all forecasts to be the value of the last observation. That is,
$$
\hat{y}_{T+h \mid T}=y_T .
$$
This method works remarkably well for many economic and financial time series.


\textbf{Drift method}

A variation on the naïve method is to allow the forecasts to increase or decrease over time, where the amount of change over time (called the drift) is set to be the average change seen in the historical data. Thus the forecast for time $T+h$ is given by
$$
\hat{y}_{T+h \mid T}=y_T+\frac{h}{T-1} \sum_{t=2}^T\left(y_t-y_{t-1}\right)=y_T+h\left(\frac{y_T-y_1}{T-1}\right) .
$$


Forecasts produced using exponential smoothing methods are weighted averages of past observations, with the weights decaying exponentially as the observations get older.

\textbf{Auto-regressive Model}

\textbf{a regression decision tree} which follows a Classification and Regression Tree (CART) algorithm [5], choosing splits to maximize the chosen split-criterion gain.  the model is relatively simple and does not take into consideration any time series aspect of the data. 








\section{Key Modeling Approaches}


\subsection{Autoregressive Approaches}

\begin{itemize}

    \item  In the first, a structural time series is treated as a state space model. In this approach, the values of the time series we observe are generated by latent space dynamics. This encompasses an enormously broad and flexible class of models, including the likes of ARIMA and the Kalman filter. Popular open source packages like bsts (Bayesian Structural Time Series, in R) and the TensorFlow Probability sts module support the state space model formulation of structural time series.
    
    \item Autoregressive models predict future values based on past values. Autoregressive models implicitly assume that the future will resemble the past.  they can prove inaccurate under certain sceanrios when system is going through a period of rapid deterioration. 

    \item  Autoregressive approaches that consider the previous few (or many) data points for each next prediction.  When applied to financial data, the ARIMA model is able to leverage the fact that financial time series data is generally related to past values. Provided there are no sudden changes in value or behavior, an ARIMA model will also be very effective for financial time series forecasting. 

\end{itemize}

\subsection{Generalized Additive Models}

\begin{itemize}

    \item In the second approach, structural time series are generalized additive models (GAMs), where the time series is decomposed into functions that add together (each of which may itself be the product or sum of multiple components) to reproduce the original series. This approach casts the time series problem as curve fitting, which does incur some trade-offs. 
    
    \item  the advantage of being scalable and easy to interpret, and tend to perform well on a large number of time series problems.

    \item Generalized additive models were formalized in their namesake paper by Hastie and Tibshirani in 1986. They replace the terms in linear regression with smooth functions of the predictors. The form of those functions determines the structure, and thus flexibility, of the model. When applied to univariate time series, the observed variable is deconstructed into smooth functions of time. There are some functions that are common to many time series, which we will outline here. Schematically, the model we are describing looks something like this:

\end{itemize}


\section{Considerations for Model Validation:}

see section 3.4 https://otexts.com/fpp2/decomposition.html
https://www.ncbi.nlm.nih.gov/pmc/articles/PMC5870978/




The fast and powerful methods that we rely on in machine learning, such as using train-test splits and k-fold cross validation, do not work in the case of time series data. This is because they ignore the temporal components inherent in the problem.

1. Train-Test split** that respect temporal order of observations.
2. Multiple Train-Test splits** that respect temporal order of observations.
3. Walk-Forward Validation** where a model may be updated each time step new data is received.


\section{Case Studies: }

http://cs229.stanford.edu/proj2017/final-reports/5244336.pdf 

http://cs229.stanford.edu/proj2017/final-reports/5244275.pdf

http://cs230.stanford.edu/projects_spring_2019/reports/18680194.pdf



- \href{https://www.nature.com/articles/s41598-022-23154-4}{Prediction of global omicron pandemic using ARIMA, MLR, and Prophet models}



\section{Predictive Monitoring Using AutoML}


The toolkit is built on top of Ray (a distributed framework for emerging AI applications open-sourced by UC Berkeley RISELab), 

- We can use Databricks AutoML using the API by doing import databricks.automl

- Foracsting examples: https://docs.databricks.com/_static/notebooks/machine-learning/automl-forecasting-example.html

- trend analysis can be done : https://www.databricks.com/blog/2022/02/09/jumpstart-your-machine-learning-forecasting-with-databricks-automl.html

- Forecasting parameters can be found here: https://docs.databricks.com/machine-learning/automl/train-ml-model-automl-api.html (e.g we need to specify target_col, HORIZON, EVALUATION METRIC)


We must provide the column that the model should predict in the `target_col` argument and the time column.  This example also specifies:

- `horizon=30` to specify that AutoML should forecast 30 days into the future. 
- `frequency="d"` to specify that a forecast should be provided for each day. 
- `primary_metric="mdape"` to specify the metric to optimize for during training.



- Hyperparameter tuning: AutoML automatically distributes hyperparameter tuning trials across the worker nodes of a cluster.
- Displays the results and provides a Python notebook with the source code for each trial run so you can review, reproduce, and modify the code. 

General critique of automl solution: https://neptune.ai/blog/automl-solutions


\section{Research Frontiers}

\href{https://towardsdatascience.com/advances-in-deep-learning-for-time-series-forecasting-and-classification-winter-2023-edition-6617c203c1d1}{Advances in Deep Learning for Time Series Forecasting and Classification}

\href{https://www.deepmind.com/blog/nowcasting-the-next-hour-of-rain}{Nowcasting the next hour of rain
}

\href{https://ai.googleblog.com/2020/12/using-automl-for-time-series-forecasting.html}{Using AutoML for Time Series Forecasting
}

\section{Resources}


- \href{https://www.pnas.org/doi/10.1073/pnas.0701020104}{On the trend, detrending, and variability of nonlinear and nonstationary time series}

- Intro to prophet: \url{https://towardsdatascience.com/time-series-analysis-with-facebook-prophet-how-it-works-and-how-to-use-it-f15ecf2c0e3a}


\url{https://machinelearningmastery.com/backtest-machine-learning-models-time-series-forecasting/}

Statistical Rethinking: 
\url{https://xcelab.net/rm/statistical-rethinking/}

Forecasting: Principles and Practice: 
\url{https://otexts.com/fpp2/what-can-be-forecast.html}

- https://structural-time-series.fastforwardlabs.com

- Practical Time Series Analysis - Prediction with statistics and Machine  Learning (O Reilley) 



\chapter{Decision Trees, XGBoost, Gradeint Boost}

- Chapter 
- Lecture 1-4 Ml for coders 
- youtube videos on the topic 


\section{Introduction To decision Trees}




\section{Research Frontiers}

- \href{https://arxiv.org/pdf/2207.08815.pdf}{Why do tree-based models still outperform deep
learning on tabular data?}

- \href{https://arxiv.org/pdf/2106.03253.pdf}{TABULAR DATA: DEEP LEARNING IS NOT ALL YOU NEED}


\section{Interviews: Key Concepts and Common Questions}


\subsection{Similarly, what does the theory of learning and then information theory tell us about generalisation ? }


\subsection{What about when data is non-idd and why does data drift occur in production?}


\subsection{How do we know if our models generalize well?}

As discussed above, In Machine Learning we are interested in finding out How well our model will perform on instances it has never seen before. We Divide the dataset into train and test set and the error rate on new cases is called generalisation error. We can get a rough estimate of it by evaluating our model on the test set. We have the option to alter the hypothesis space (i.e the set of functions that the learning algorithm is allowed to select as being the solution). Our challenge here is to choose a model with the right 'capacity'. This means we pick Machine Learning models that match the complexity of task to be solved and the training data available. e.g quadratic model will generalize well if the underlying function is quadratic.

To better understand generalization, we can decompose the generalization error into bias and variance.

\textbf{High Variance:} When a learning algorithm outputs a prediction function g with generalization error much higher than its training error(i.e gap between the training error and test error is too large.), then it is said to have overfit the training data[2].

\textbf{High Bias:}  Means the learning algorithm is underfitting the data. In practice we can look at the train set error and dev set error. If its not fitting the training set well (as shown by the train set error) then algorithm is said to have high bias. An algorithm can have both as well if the train set error is high and gap between train set and test set is high as well.

Overall, we have to draw a balance between low training error while at the same time try to achieve a small gap between training and test error(generalization). We keep in mind that measuring generalization error this way is proxy method as we don't know the ground truth 'f' nor the distribution 'D'(that we talked about earlier).

\textbf{no free lunch” theorem:}
According to “no free lunch” theorem, no learner can beat random guessing over all possible functions to be learned.

A \textit{model’s capacity} measures the range of functions it can fit. VC dimension measures the capacity of a binary classifier.  \textit{Occam’s Razor} states that, with the same expressiveness, simple models are preferred.  Training error and generalization error versus model capacity usually form a \textit{U-shape relationship}. We find the optimal capacity to achieve low training error and small gap between training error and generalization error.  Bias measures the expected deviation of the estimator from thetrue value; while variance measures the deviation of the estimator from the expected value, or variance of the estimator.  As model capacity increases, bias tends to decrease, while variance tends to increase, yielding another U-shape relationship between generalization error versus model capacity. \textit{We try to find the optimal capacity point}, of which under-fitting occurs on the left and over-fitting occurs on the right. \textit{ Regularization } add a penalty term to the cost function, to reduce the general-ization error, but not training error.  \textit{No free lunch theorem} states that there is no universally bestmodel, or best regularizor. An implication is that deep learning may not be the best model for some problems. There are model parameters, and hyperparameters for model capacity and regularization. \textit{Cross-validation} is used to tune hyperparameters, to strike a balance between bias and variance, and to select the optimal model. 





\section{Few common concepts}


\begin{itemize}


    \item \textbf{Cost Function:} It allows us to see how bad the model is performing. We can measure distance between model's predictions and data provided by the examples



    \item \textbf{Model Parameters:} In linear regression they are theta 1 and theta2 We find parameters that optimally fit best to the data. 
    
    \item \textbf{Training: }During training we tweak these parameters allowing us fit the training set and make good predictions on the the new cases 

    \item \textbf{Data Quantity}: inussificent quanity of data can be problametic. Simple methods can beat advanced ones if sufficient data is available. Overall, its crucial to use a training set that is representative of the cases youi want to generalize to. If the data is too small, you will have sampling noise. i.e non-representative data as a result of chance

    \item \textbf{Sampling bias} Even very large samples can be non-representative if sampling method is flawed. This is called sampling bias. 

    \item \textbf{Feature Engineering} It is important that come up with a good set of features to train on. This involves Feature selection, feature extraction (combining existing features to produce a more useful one or using dimensionality reduction). Lastly creating new features by gathering new data. 

    \item \textbf{Overfitting the training data:} In real life, we can see one example and make conclusions about that. In ML, overfitting is when model performs well on the training data but doens't generalize well. It happens when model is too complex relative to the amount and nouiseiness of the training data. We can gather more data or reduce noise from the data or simplify the model with fewer parameters. Constraining the to make it simpler and reduce the risk of overfitting is called regularization. e.g if linear model has two parameters then it has two degrees of freedom to adapt to the model to the training data. In ML it comes down to fitting the model to the data perfectly and yet keeping it simple enough so it can generalize. This can be controlled by hyperparemeters. If the model is too simple it might underfit and if its too complex it will overfit. 

    \item \textbf{Testing and validation} We can hold out 20\% of the data for testing (depending on size). The error rate on new cases is called generalization error. We can get an estimate of this by evaluating our model on the test set. If the training error is low bu the generalization error is high, it means that your model is overfitting the training data.  

    \item \textbf{Hold-out validation:} Repeatedly evaluating a model on test set can be problematic when selecting multiple candidate models. A common solution to this problem is called hold-out validation: you can simply hold out a part of the training set to evaluate several candaidate models and select the best one. The new held out set is called the validation set. More specifically, you can training multiple models with various hyperparameters on the reduced training set and you select the model that performs best on the valluidation set. After theis holdout valudation process you train the best model on the training set and this fives you the final model. Lastly, you evaluate this final model on the test set to get an estimate of the generalization error. 
    
    \item \textbf{Cross-validation} If validation set is too small, then model evaluations will be imprecise. Cross-validation involves many small validation sets. each model is evaluated once per validation set after it is trained on the rest of the data. By averaging out all the evaluations of a model, you get a much more accurate measure of its performance. 

    \item \textbf{Scikit Learn's K-fold cross validation} We can randomly split the training set into 10 distict subsets called folds, then it trains and evaluates the decision tree model 10 times, picking a different fold for evaluation every time and training on other 9 folds. 

    \item \textbf{No free lunch theorem} without making any assumptions avout the data there is no apriori way of knowing if one model is going to perform better over the other. 
 
\end{itemize}















\textbf{How do we know if our models generalizes well?}

As discussed above, In Machine Learning we are interested in finding out How well our model will perform on instances it has never seen before. We Divide the dataset into train and test set and the error rate on new cases is called generalisation error. We can get a rough estimate of it by evaluating our model on the test set. We have the option to alter the hypothesis space (i.e the set of functions that the learning algorithm is allowed to select as being the solution). Our challenge here is to choose a model with the right 'capacity'. This means we pick Machine Learning models that match the complexity of task to be solved and the training data available. e.g quadratic model will generalize well if the underlying function is quadratic.

To better understand generalization, we can decompose the generalization error into bias and variance.

\textbf{High Variance:} When a learning algorithm outputs a prediction function g with generalization error much higher than its training error(i.e gap between the training error and test error is too large.), then it is said to have overfit the training data[2].

\textbf{High Bias:}  Means the learning algorithm is underfitting the data. In practice we can look at the train set error and dev set error. If its not fitting the training set well (as shown by the train set error) then algorithm is said to have high bias. An algorithm can have both as well if the train set error is high and gap between train set and test set is high as well.

Overall, we have to draw a balance between low training error while at the same time try to achieve a small gap between training and test error(generalization). We keep in mind that measuring generalization error this way is proxy method as we don't know the ground truth 'f' nor the distribution 'D'(that we talked about earlier).


\textbf{Wisdom nugget: Model Capacity}

In particular, one wants to avoid learning a prediction function that effectively memorizes the training sample — this would be “overfitting” to the particular data at hand. 

%We’ll introduce you to a clean and useful conceptual framework, rooted in statistical learning theory, for reasoning about these issues.
We’ll also cover useful techniques for avoid overfitting in the practically common case where the amount of training data is relatively small. One such technique is “regularization,” whereby you penalize “more complex” prediction functions in favor of “simpler” ones. This is a version of Occam’s Razor, which advocates accepting the simplest theory that explains the known observations.


A \textit{model’s capacity} measures the range of functions it can fit. VC dimension measures the capacity of a binary classifier.  \textit{Occam’s Razor} states that, with the same expressiveness, simple models are preferred.  Training error and generalization error versus model capacity usually form a \textit{U-shape relationship}. We find the optimal capacity to achieve low training error and small gap between training error and generalization error.  Bias measures the expected deviation of the estimator from thetrue value; while variance measures the deviation of the estimator from the expected value, or variance of the estimator.  As model capacity increases, bias tends to decrease, while variance tends to increase, yielding another U-shape relationship between generalization error versus model capacity. \textit{We try to find the optimal capacity point}, of which under-fitting occurs on the left and over-fitting occurs on the right. \textit{ Regularization } add a penalty term to the cost function, to reduce the general-ization error, but not training error.  \textit{No free lunch theorem} states that there is no universally bestmodel, or best regularizor. An implication is that deep learning may not be the best model for some problems. There are model parameters, and hyperparameters for model capacity and regularization. \textit{Cross-validation} is used to tune hyperparameters, to strike a balance between bias and variance, and to select the optimal model. 

see: 
https://www.overleaf.com/project/60b980221babf728106d2604 



Withoput making assumptions about the data, there is no way of knowing which is the right model. So we perform experiments. Our goal in experiments is to find the model with the right capacity and one that can generalize beyond the seen examples. if the model is too simple, it will underfit, if its too complex, it will overfit. In practice we can find out the generalization error by evaluating it on test set. 



\section{Law of Large Numbers}

 As the number of trials or observations increases, the actual or observed probability approaches the theoretical or expected probability.




\chapter{Support Vector Machines}



\chapter{Recommender Systems}


\section{Recommender System System design}

https://medium.com/double-pointer/system-design-interview-recommendation-system-design-as-used-by-youtube-netflix-etc-c457aaec3ab


\section{Evaluating Recommender Systems}


\textbf{Normalized Discounted Cumulative Gain (NDCG):} normalized sum of discounted relevance values of all the results

\textbf{Hit Rate (HR, HR@N):} how often does the list of recommendations include something that's actually relevant?

\textbf{Mean Reciprocal Rank (MRR):} where is the first relevant item in the list of recommendations?

Coverage: what's the percentage of items seen in raining data that are also seen in recommendations?

RecList: behavioral testing for RecSys

Paper: Beyond NDCG: behavioral testing of recommender systems with RecList

article: 

- https://towardsdatascience.com/ndcg-is-not-all-you-need-24eb6d2f1227
- https://towardsdatascience.com/ndcg-is-not-all-you-need-24eb6d2f1227


library: 
code: https://colab.research.google.com/drive/1Wn5mm0csEkyWqmBBDxNBkfGR6CNfWeH-?usp=sharing



\section{Research}


\href{https://eugeneyan.com/writing/recsys2022/?utm_campaign=Data_Elixir&utm_source=Data_Elixir_408/}{RecSys 2022: Recap, Favorite Papers, and Lessons
}



\chapter{Unsupervised Learning}


\section{Dimensionality reduction}
One could also call this topic “the unreasonable effectiveness of sophomore-level linear algebra.” This is a major topic, and it will
occupy us for three weeks. Many data sets are usefully interpreted as points in space
(and hence matrices, with the vectors forming the rows or the columns of a matrix).
For example, a document can be mapped to a vector of word frequencies. Graphs
(social networks, etc.) can also usefully be viewed as matrices in various ways. We’ll see that linear algebraic methods are incredibly useful for exposing the “geometry”
of a data set, and this allows one to see patterns in the data that would be otherwise undetectable. Exhibit A is principle component analysis (PCA), which identifies
the “most meaningful” dimensions of a data set. We’ll also cover the singular value
decomposition (SVD), which identifies low-rank structure and is useful for denoising
data and recovering missing data. Finally, we’ll see how eigenvalues and eigenvectors
have shockingly meaningful interpretations in network data. Linear algebra is a muchmaligned topic in computer science circles, but hopefully the geometric intuition and
real-world applications we provide will bring the subject to life.


Why is the curse of dimensionality a problem? The issue is that the natural representation of data is often high-dimensional. Recall our motivating examples for representing data points in real space. With documents and the bag-of-words model, the number of coordinates equals the number of words in the dictionary — often in the tens of thousands. Images are often represented as real vectors, with at least one dimension per pixel (recording pixel intensities) — here again, the number of dimensions is typically in the tens of thousands. The same holds for points representing the purchase or page view history of an Amazon customer, or of the movies watched by a Netflix subscriber.

The friction between the large number of dimensions we want to use to represent data and the small number of dimensions required for computational tractability motivates dimensionality reduction. The goal is to re-represent points in high-dimensional space as points in low-dimensional space, preserving interpoint distances as much as possible. We can think of dimensionality reduction as a form of lossy compression, tailored to approximately preserve distances. 


\section{The nearest neighbor problem}

Dimension reduction continues the theme of lossy compression:  it's about compressing data while approximately preserving similarity information (represented using distances). In the nearest neighbor  problem, you are given a point  set  (e.g.,  representing  documents) and want topreprocess it so that, given a query (e.g., representing a keyword search query), you can quickly determine which point is closest to the query.  This problem oers our rstmethod of understanding and exploring a data set




\section{Principle Component Analysis}

Assume that your grandma likes wine and would like to find characteristics that best describe wine bottles sitting in her cellar. There are many characteristics we can use to describe a bottle of wine including age, price, color, alcoholic content, sweetness, acidity, etc. Many of these characteristics are related and therefore redundant. Is there a way we can choose fewer characteristics to describe our wine and answer questions such as: which two bottles of wine differ the most?

PCA is a technique to construct new characteristics out of the existing characteristics. For example, a new characteristic might be computed as age - acidity + price or something like that, which we call a linear combination.

To differentiate our wines, we'd like to find characteristics that strongly differ across wines. If we find a new characteristic that is the same for most of the wines, then it wouldn't be very useful. PCA looks for characteristics that show as much variation across wines as possible, out of all linear combinations of existing characteristics. These constructed characteristics are principal components of our wines.

If you want to see a more detailed, intuitive explanation of PCA with visualization, check out amoeba's answer on StackOverflow. This is possibly the best PCA explanation I've ever read.

https://stats.stackexchange.com/questions/2691/making-sense-of-principal-component-analysis-eigenvectors-eigenvalues/140579#140579 

- Can be used for dimensionality reduction
- Rotate our axis such they point in the direction in which data varies the most
- That way we can capture as much information as possible in a smaller number of dimensions 
- What straight line drawn through the origin explains the maximum amount of variance in my data
- What straight line orthogonal to the first principle component explains the variance thats left 


https://youtu.be/AS0H43hMoWM

\subsection{T-Distributed Stochastic Neighbor Embedding (t-SNE)}

- https://twitter.com/GCLinderman/status/1058377443835146240?s=20&t=kA2VI4pC9Tx7OdpuIK4T3A

- https://distill.pub/2016/misread-tsne/




\subsection{Unsupervised Learning}

Unsupervised algorithms experience a dataset containing many features, then learn useful properties of the structure of this dataset experience a dataset containing features,but each example is also associated with a label or target.


\subsubsection{Clustering}

-  Clustering, consists of dividing the dataset into clusters of similar examples.

- Given samples $\chi^{(1)}, \ldots, x^{(n)} \in \mathbb{R}^{D}$, the goal is to find a partitioning (or "clustering") of the samples that groups together samples that are similar. There are many different objectives, depending on the definition of the similarity between samples and exactly what criterion is to be used (e.g., minimize the average distance between elements inside a cluster and maximize the average distance between elements across clusters). Other methods perform a "soft" clustering, in which samples may be assigned $0.9$ membership in one cluster and $0.1$ in another. Clustering is sometimes used as a step in density estimation, and sometimes to find useful structure in data.

\subsubsection{Density Estimation}

- Given samples $x^{(1)}, \ldots, x^{(n)} \in \mathbb{R}^{D}$ drawn IID from some distribution $\operatorname{Pr}(X)$, the goal is to predict the probability $\operatorname{Pr}\left(x^{(n+1)}\right)$ of an element drawn from the same distribution. Density estimation sometimes plays a role as a "subroutine" in the overall learning method for supervised learning, as well.

-  In the context of deep learning, we usually want to learn the entire probability distribution that generated a dataset, whether explicitly, as in density estimation, or implicitly, for tasks like synthesis or denoising. 



\section{Interview Questions:}

\textbf{Why do we need dimensionality reduction?}

\textbf{ Eigendecomposition is a common factorization technique used for dimensionality reduction. Is the eigendecomposition of a matrix always unique?}


\textbf{ Name some applications of eigenvalues and eigenvectors}

\textbf{We want to do PCA on a dataset of multiple features in different ranges. For example, one is in the range 0-1 and one is in the range 10 - 1000. Will PCA work on this dataset?}

\textbf{ Under what conditions can one apply eigendecomposition? What about SVD?}

\textbf{What is the relationship between SVD and eigendecomposition?}

\textbf{What’s the relationship between PCA and SVD?}

\textbf{How does t-SNE (T-distributed Stochastic Neighbor Embedding) work? Why do we need it?}


\chapter{Summary of Key Steps}


%\section{General Steps of Solving an ML problem} 


\begin{comment}

Examples of System Design/Data Engineering tasks include:

\begin{itemize}

    \item \textbf{Ingest data from a data source} A dataset is a collection of many examples, as defined in section. Sometimes we call examples data points. 
    \item Build and maintain a data warehouse
    \item create a data pipeline
    \item create an analytics table for a specific use case
    \item migrate data to cloud
    \item schedule and automate pipelines
    \item backfill data
    \item debug data quality issues
    \item optimize queries
    \item design a database
    \item \textbf{Feature Representation:} Data is characterised by features. e.g take a song and turn it into a vector. 
    
\end{itemize}

\end{comment}



%\section{ML Process Summary}



Every organisation is unique and they have their own set of unique challenges.  Staying ahead in the accelerating AI race requires executives to make nimble, informed decisions about where and how to employ AI in their business. Here Alignment with strategic goals of your organization is key, where its good to Define a prioritization mechanism for the efforts. 

Generally speaking when an organisation takes an initiative and makes efforts to incorporate ML and DS into their processes, their high level goal is to for example, transform the organisation where (i) you want decisions that are being made at operational and managerial level to be supported by data and analytics. (ii) Organisations also try to leverage ML technologies to support the process of improving the efficiency and enabling automation where possible. 

To start with, we recommend designing a organisational road-map for AI and data science where we identify the opportunities for projects that are low hanging fruit. This means executing projects that have the potential for having a maximum impact in the near future.  All of this is hard work, it requires  dedicated efforts, where you often start with placing the right data infrastructure in place that can support ML or data science projects. The teams also make an effort to carefully collect the right data, analyze it manually, Create experiments that building models, or generate insights by for example, utilising intelligent algorithms and so on.  To get started, it might be good idea to develop certain pilots(prototypes) to gain trust/interest of the leardership/investors/stakeholders and then build up momentum gradually. i.e Choose a use-case easy to explain, possibly solvable with simple, interpretable AI while also Identifying business metrics to express value. 

To get started, ask the following fundamental questions: 

\begin{itemize}
    \item What is the problem we are solving?
    \item How do we know when we've won?
    \item What are we going to do with this?
    \item What are the product, user, and infrastructure concerns for the deployment of this?
\end{itemize}



The ML project development life cycle An ML project starts with data about a service or product on which you fit models to be deployed into production. These models need to be monitored and their performance evaluated. The firm’s AI infrastructure supports these tasks.

One crucial aspect of machine learning approaches to solving problems is that human engineering plays an important role. A human still has to frame the problem: acquire and organize data, design a space of possible solutions, select a learning algorithm and its parameters, apply the algorithm to the data, validate the resulting solution to decide whether it's good enough to use, etc. These steps are of great importance.

In this guide, we will discuss some of the details on approaching an ML project. Details of these steps is as follows: 

\subsection{Frame the Problem and Look at the Big Picture} \\

\begin{enumerate}
    \item identify areas of the business that could benefit from machine learning 
    \item Define the objective in business terms. 
    
    \begin{description}
     \item - appropriately frame and scope the task
     \item - communicate with other stakeholders about what machine learning is and is not capable of (there are often many misconceptions)
     \end{description}
    \item How will your solution be used?
    \begin{description}
        \item - develop understanding of business strategy, risks, and goals to make sure everyone is on the same page
     \end{description}
    \item  What are the current solutions/workarounds (if any)? (How are being things currently done)
    
    \item How should you frame this problem (supervised/unsupervised, online/offline, etc.)?
    
    \item How should performance be measured?
        \begin{description}
            \item Choosing F1-score to evaluate a model’s performance on a classification task 
            \item Implementing evaluation metrics such as accuracy, precision, recall, intersection over union, or mean average precision (mAP)
        
        \end{description}
    
    \item Is the performance measure aligned with the business objective? 
    
    \item What would be the minimum performance needed to reach the business objective?
    
    \item What are comparable problems? Can you reuse experience or tools?
    
    \item Is human expertise available?
    
    \item How would you solve the problem manually?
    
    \item  List the assumptions you (or others) have made so far. (and  Verify assumptions if possible.)

\end{enumerate}


\subsection{Data Acquizition} 

\textit{Note: automate as much as possible so you can easily get fresh data.
}

\begin{enumerate}

    \item List the data you need and how much you need.
    \item  Find and document where you can get that data (identify what kind of data the organization has). e. g Querying data from a database
    \item Check how much space it take and engineering efforts required(initial assessment). 
    
    \item  understand operational constraints (e.g. what data is actually available at inference time)
    
    \item  Legal/Ethics/Privacy
    
    \begin{description}
        \item  Get access authorizations and check legal obligations 
        \item proactively identify ethical risks, including how your work could be mis-used by harassers, trolls, authoritarian governments, or for propaganda/disinformation campaigns (and plan how to reduce these risks)
        \item  identify [potential biases and potential negative feedback loops]
        \item Ensure sensitive Information is deleted or protected (e.g anonymized)
        \item  Keep track of data sources. e.g  Setting up a data version control system
    \end{description}
    
\end{enumerate}



\textbf{Examples: }

\begin{itemize}
    \item Setting up a Mechanical Turk project 
    \item Collecting data by manually taking images of cats 
    \item Coding a JavaScript tracker on a website to collect user data 
    \item Scraping the web and, if necessary, synchronizing data from different sources
    \item  Labeling data. e.g \- Writing a labeling tutorial for workers
\end{itemize} 



\subsection{EDA - Explore the data:}

\textbf{Before EDA:}

\begin{itemize}
    \item Check the size and type of data  
    \item See if the data is in appropriate for  - Convert the data to a format you can easily manupulate (without changing the data itself)
    \item  Sample a test set, set it aside and never look at it
\end{itemize}



\textbf{EDA:}


\begin{enumerate}
    \item Grab a copy of the data 
    \item  Document the EDA in a Jupyter notebook
    \item  Study Each Attribute and its characteristics (Name, Type, \% of missing values, noisy, usefulness, type of distribution)
    \item For Supervised Machine Learning, identify the target attribute 
    \item  Visualize the data. e.g Visualizing high-dimensional data in lower dimensions using methods such as PCA or t-SNE
    \item  Study the correlations between attributes 
    \item Study how you would solve the problem manually
    \item  Identify the promising transformations you may want to apply
    \item make plans to collect more of different data (if needed and if possible)
\end{enumerate}

\subsection{Data Pre-Processing}


\textbf{Data Cleaning:} 

\begin{itemize}
    \item Fix or remove outliers 
    \item  Fill in missing values(e.g with zero, mean, median) or drop their rows  
    \item deal with corrupted/erronous  data
    \item (For ML) Feature Selection: Drop the attributes that provide no useful infomration for the task 
\end{itemize}

\textbf{Prepare the data:}

\begin{itemize}
    \item Create appropriate training and validation set https://www.fast.ai/2017/11/13/validation-sets/)
    \item Writing function for data transformation (they will come in handy next time) e.g  Converting a continuous feature into a categorical feature using bucketing 
    \item  stitch together data from many different sources: often this data has been collected in different formats or with inconsistent conventions
    
    \item  Moving data and building data pipelines. e.g     \begin{description}
            \item - Writing a script to allow online learning for a model 
            \item - Designing an ETL system 
            \item - Writing a script to pre-process training data and send it as input to a model automatically
            \item - Writing a script to record model predictions in a database
        \end{description}
\end{itemize}


\textbf{ Feature Engineering: }

\begin{itemize}
    \item Discretize continuous features
    \item Decompose features (e.g categorical, data/time etc)
    \item Add Promising transformations of feature(e.g $log(x), sort(x), x^2$)
    \item Aggregate features into promising new features
    \item Standardise or normalize features  
\end{itemize}


\subsection{Modeling:}

\begin{enumerate}

    \item Train many quick and dirty models from different categories using standard parameters
    \item Measure and compare their performance 
    \item  Analyse the most significant variables for each algorithm
    \item Analyse the types of errors the models make 
    \item Perform a quick round of feature selection and engineering
    \item  Perform one or two more quick iterations of the five previous steps 
    \item Shortlist the top three to five most promising models, preferrring models that make different types of errors 
    \item fit model resource needs into constraints (e.g. will the completed model need to run on an edge device, in a low memory or high latency environment, etc)

\end{enumerate}

\subsection{Fine-tunings} 

\begin{itemize}
    \item Fine-tune the hyper hyperparameters using cross-validation (e.g. in the case of deep learning, this includes choosing an architecture, loss function, and optimizer) - *ADD MORE DETAILS HERE* 
    
    train the model (and debug why it’s not training). This can involve:
     
    \item adjusting hyperparmeters (e.g. such as the learning rate)
    
    \item outputing intermediate results to see how the loss, training error, and validation error are changing with time
    
    \item inspecting the data the model is wrong on to look for patterns
    
    \item identifying underlying errors or issues with the data
    
    \item realizing you need to change how you clean and pre-process the data
    
\end{itemize}


  - 
  - realizing you need more or different data augmentation
  - realizing you need more or different data
  - trying out different models
  - identifying if you are under- or over-fitting

* try ensemble methods. Combining your best models will often produce better performance than running them individually. 
* Once you are confident about your final model, measure its performance on the test set to esitmate the generalization error. 


\subsection{Present/Launch the solution}

\textbf{Productionize} \\

\begin{itemize}
    \item creating an API or web app with your model as an endpoint in order to productionize
    \item exporting your model into the needed format
    \item plan for how often your model will need to be retrained with updated data (e.g. perhaps you will retrain nightly or weekly)
\end{itemize}

\textbf{Monitor} \\ 

\begin{itemize}
    \item track model performance over time
    \item monitor the input data, to identify if it changes with time in a way that would invalidate your model
    \item communicate your results to the rest of the organization
    \item have a plan in place for how you will monitor and respond to mistakes or unexpected consequences
\end{itemize}





\chapter{Model Deployment and Monitoring}



\section{Resources:}

file:///Users/asjad/Downloads/Machine%20Learning.pdf
https://arxiv.org/pdf/2206.13446.pdf

- Kevin Murphy's Book 
- https://mlu-explain.github.io/ 
- https://aws.amazon.com/machine-learning/mlu/


\chapter{Case Studies}


Case studies

To learn to design machine learning systems, it's helpful to read case studies to see how actual teams deal with different deployment requirements and constraints. Many companies, such as Airbnb, Lyft, Uber, Netflix, run excellent tech blogs where they share their experience using machine learning to improve their products and/or processes. If you're interested in a company, you should visit their tech blogs to see what they've been working on -- it might come up during your interviews! Below are some of these great case studies.

    Using Machine Learning to Predict Value of Homes On Airbnb (Robert Chang, Airbnb Engineering & Data Science, 2017)

    In this detailed and well-written blog post, Chang described how Airbnb used machine learning to predict an important business metric: the value of homes on Airbnb. It walks you through the entire workflow: feature engineering, model selection, prototyping, moving prototypes to production. It's completed with lessons learned, tools used, and code snippets too.

    Using Machine Learning to Improve Streaming Quality at Netflix (Chaitanya Ekanadham, Netflix Technology Blog, 2018)

    As of 2018, Netflix streams to over 117M members worldwide, half of those living outside the US. This blog post describes some of their technical challenges and how they use machine learning to overcome these challenges, including to predict the network quality, detect device anomaly, and allocate resources for predictive caching.

    150 Successful Machine Learning Models: 6 Lessons Learned at Booking.com (Bernardi et al., KDD, 2019).

    As of 2019, Booking.com has around 150 machine learning models in production. These models solve a wide range of prediction (e.g. predicting users' travel preferences and how many people they travel with) and optimization (e.g.optimizing the background images and reviews to show for each user). Adrian Colyer gave a good summary of the six lessons learned here:
        Machine learned models deliver strong business value.
        Model performance is not the same as business performance.
        Be clear about the problem you're trying to solve.
        Prediction serving latency matters.
        Get early feedback on model quality.
        Test the business impact of your models using randomized controlled trials.

    How we grew from 0 to 4 million women on our fashion app, with a vertical machine learning approach (Gabriel Aldamiz, HackerNoon, 2018)

    To offer automated outfit advice, Chicisimo tried to qualify people's fashion taste using machine learning. Due to the ambiguous nature of the task, the biggest challenges are framing the problem and collecting the data for it, both challenges are addressed by the article. It also covers the problem that every consumer app struggles with: user retention.

    Machine Learning-Powered Search Ranking of Airbnb Experiences (Mihajlo Grbovic, Airbnb Engineering & Data Science, 2019)

    This article walks you step by step through a canonical example of the ranking and recommendation problem. Four main steps are system design, personalization, online scoring, and business aspect. The article explains which features to use, how to collect data and label it, why they chose Gradient Boosted Decision Tree, which testing metrics to use, what heuristics to take into account while ranking results, how to do A/B testing during deployment. Another wonderful thing about this post is that it also covers personalization to rank results differently for different users.

    From shallow to deep learning in fraud (Hao Yi Ong, Lyft Engineering, 2018)

    Fraud detection is one of the earliest use cases of machine learning in industry. This article explores the evolution of fraud detection algorithms used at Lyft. At first, an algorithm as simple as logistic regression with engineered features was enough to catch most fraud cases. Its simplicity allowed the team to understand the importance of different features. Later, when fraud techniques have become too sophisticated, more complex models are required. This article explores the tradeoff between complexity and interpretability, performance and ease of deployment.

    Space, Time and Groceries (Jeremy Stanley, Tech at Instacart, 2017)

    Instacart uses machine learning to solve the task of path optimization: how to most efficiently assign tasks for multiple shoppers and find the optimal paths for them. The article explains the entire process of system design, from framing the problem, collecting data, algorithm and metric selection, topped with tutorial for beautiful visualization.

    Uber's Big Data Platform: 100+ Petabytes with Minute Latency (Reza Shiftehfar, Uber Engineering, 2018)

    With massive data comes massive engineering requirement. Relying heavily on data for decision making, "from forecasting rider demand during high traffic events to identifying and addressing bottlenecks in our driver-partner sign-up process", Uber has collected "over 100 petabytes of data that needs to be cleaned, stored, and served with minimum latency." This article focuses on the evolution of analytical data warehouse at Uber, from Vertica to Hadoop to their own Spark library Hudi, each with their limitations analyzed and addressed.

    Creating a Modern OCR Pipeline Using Computer Vision and Deep Learning (Brad Neuberg, Dropbox Engineering, 2017)

    An application as simple as a document scanner has two distinct components: optical character recognition and word detector. Each requires their own production pipeline, and the end-to-end system requires additional steps for training and tuning. This article also goes into detail the team's effort to collect data, which includes building their own data annotation platform.

    Scaling Machine Learning at Uber with Michelangelo (Jeremy Hermann and Mike Del Balso, Uber Engineering, 2019)

    Uber uses extensive machine learning in their production, and this article gives an impressive overview of their end-to-end workflow, where machine learning is being applied at Uber, and how their teams are organized.

    Deep Learning for Recommender Systems (Justin Basilico, Research/Engineering at Netflix, 2018)

    Recommendations generate over $1B yearly revenue for Netflix. The joy of spending few seconds to find something great to watch, directly impacts customer satisfaction.

    Making Netflix Machine Learning Algorithms Reliable (Justin Basilico, Research/Engineering at Netflix, 2017)

    Netflix develops a variety of machine learning algorithms (including regression, factorization, topic modeling, ensemble learning, neural networks, bandits, etc) to a variety of problems (personalized ranking, trending now, video similarities, search, top-n ranking, etc), to help members find content to watch and enjoy.




\chapter{Useful Books and Courses for Furhter learnig}

\section{About this document and Learning Roadmap}


what are the key topics? and for each topic what is the useful set of bare minimum theory that an ML practitioner should know? Here are a few learning roadmaps 

\begin{itemize}
    \item \href{https://i.am.ai/roadmap/#note }{AI Roadmap} 
    \item \href{https://github.com/dformoso/machine-learning-mindmap
}{Machine Learning MindMap} 
    \item \href{https://madewithml.com/}{Made with ML} 
    \item \href{https://whimsical.com/machine-learning-roadmap-2020-CA7f3ykvXpnJ9Az32vYXva }{Machine Learning Roadmap}
    \item \href{https://www.youtube.com/watch?v=vcE9WGbi4QY
}{Cornell Applied Machine Learning Course} 

\end{itemize}



\section{ML foundations}

- Andrew NG: https://irosyadi.gitbook.io/irosyadi/course/machine-learning-andrewng


\section{Interview}

ML Basics and ML Systems: 

- https://huyenchip.com/ml-interviews-book/
- https://mlengineer.io/machine-learning-interview-books-37dbbb09416d
- Deep Learning interviews 


\section{Courses:}

- Corenell


\end{document}
